% 极简稳定版：信号分析与处理课程笔记主文件
\documentclass[oneside, UTF8]{ctexbook}

% ========== 核心宏包（补齐缺失项+优化） ==========
% 数学公式基础（核心）
\usepackage{amsmath,amssymb,mathtools}
% 跨页公式/编号优化
\usepackage{empheq}
% 图片/图形支持
\usepackage{graphicx}
% 页面布局
\usepackage{geometry}
% 表格美化（\toprule/\midrule/\bottomrule）
\usepackage{booktabs}
% 长表格支持
\usepackage{longtable}
% 列表环境优化
\usepackage{enumitem}
% TikZ绘图核心（含必备库）
\usepackage{tikz}
\usetikzlibrary{
    positioning,    % 节点定位（解决of语法报错）
    arrows,         % 箭头样式优化
    shapes,         % 基础图形样式
    calc            % 坐标计算
}
% 信号绘图专用：PGFPlots
\usepackage{pgfplots}
\pgfplotsset{
    compat=1.17,    % 指定兼容版本，避免警告
    width=0.8\textwidth, % 绘图宽度适配页面
    height=6cm      % 统一绘图高度
}
% 超链接/URL（最后加载hyperref，避免选项冲突）
\usepackage[
    colorlinks=true,        % 彩色链接（而非边框）
    linkcolor=black,        % 内部链接颜色（目录/章节）
    citecolor=blue,         % 引用颜色
    urlcolor=blue,          % URL颜色
    pdfstartview={FitH}     % PDF默认视图
]{hyperref}
\usepackage{url} % url包放在hyperref后，避免冲突

% ========== 全局命令定义（解决子文件未定义问题） ==========
% 抽样函数（子文件\Sa）- 用providecommand避免重复定义
\providecommand{\Sa}{\operatorname{sinc}}
% 冲激函数（统一格式）
\providecommand{\delt}{\delta}
% 阶跃函数（统一格式）
\providecommand{\ufunc}{u}
% 斜坡函数（统一格式）
\providecommand{\rfunc}{r}

% ========== 页面与布局配置（优化中文排版） ==========
\geometry{
    a4paper, 
    margin=2.5cm,
    headsep=10pt,   % 页眉间距
    footskip=15pt   % 页脚间距
}
% 行距优化（中文阅读更舒适）
\linespread{1.2}
% 段落间距优化
\setlength{\parskip}{0.5em}
\setlength{\parindent}{2em} % 中文首行缩进2字符

% ========== 屏蔽无关警告（最大化整洁度） ==========
\usepackage{silence}
% 屏蔽xeCJK字体重定义警告
\WarningFilter{xeCJK}{Redefining CJKfamily}
% 屏蔽字体形状警告
\WarningFilter{LaTeX Font}{Some font shapes were not available}
\WarningFilter{LaTeX Font}{Font shape `TU/'}
% 屏蔽pgfplots版本兼容警告
\WarningFilter{pgfplots}{compat=1.17}
% 屏蔽长表格轻微排版警告
\WarningFilter{longtable}{Underfull \hbox}
\WarningFilter{longtable}{Overfull \hbox}
% 屏蔽hyperref所有警告（包括bookmarks重复）
\WarningFilter{hyperref}{Option `bookmarks' has already been used}
\WarningFilter{hyperref}{Token not allowed in a PDF string}
% 屏蔽cSpell拼写检查警告（LaTeX专用词汇）
\WarningFilter{cSpell}{Unknown word}
% 全局屏蔽轻微排版警告
\hbadness=10000
\vfuzz=100pt
\vbadness=10000

% ========== 基础配置（目录/标题） ==========
% 目录/参考文献绑定
\usepackage{tocbibind}
\setcounter{tocdepth}{3}    % 目录显示到三级标题
\setcounter{secnumdepth}{3} % 标题编号到三级标题

% 标题格式优化（适配中文）
\usepackage{titlesec}
\titleformat{\chapter}{\bfseries\Large}{\thechapter}{1em}{}
\titleformat{\section}{\bfseries\large}{\thesection}{1em}{}
\titleformat{\subsection}{\normalsize\bfseries}{\thesubsection}{1em}{}
\titleformat{\subsubsection}{\normalsize}{\thesubsubsection}{1em}{}
% 标题间距优化
\titlespacing{\chapter}{0pt}{20pt}{10pt}
\titlespacing{\section}{0pt}{15pt}{5pt}
\titlespacing{\subsection}{0pt}{10pt}{3pt}

% ========== TikZ/PGFPlots通用样式（信号绘图专用） ==========
\tikzset{
    % 基础信号绘图样式
    signal plot/.style={
        thick,
        scale=1.0,
        font=\small,
        every axis plot/.append style={ultra thick, blue},
        every axis/.append style={
            axis lines = middle,    % 坐标轴居中
            xlabel={$t$},           % 时域默认x轴标签
            ylabel={$f(t)$},        % 时域默认y轴标签
            grid = major,           % 主网格
            grid style = {dashed, gray!30}, % 网格样式
            tick label style={font=\small},
            label style={font=\small}
        }
    },
    % 信号关系图样式（如微积分关系）
    signal relation/.style={
        >=stealth,        % 箭头样式
        node distance=1.5cm,
        font=\small,
        every node/.style={align=center}
    }
}

% ========== 超链接补充配置（避免重复设置bookmarks） ==========
\hypersetup{
    pdfauthor={2024级电气工程及其自动化专业}, % PDF作者信息
    pdftitle={信号分析与处理（AUTO2005）课程笔记} % PDF标题
}

% ========== 元信息 ==========
\title{信号分析与处理（AUTO2005）课程笔记}
\author{2024级电气工程及其自动化专业}
\date{2026春季学期}

\begin{document}
% 封面（缺失时注释，不影响编译）
% % 信号分析与处理专用封面（纯LaTeX原生语法，无额外依赖）
% 方案：先删除已有命令（若存在），再重新定义，避免重复
\providecommand{\makeSignalsCover}{} % 先声明命令（避免未定义）
\renewcommand{\makeSignalsCover}{
    \begin{titlepage}
      \newgeometry{left=2cm, right=2cm, top=2cm, bottom=2cm}
      \centering
      \vspace*{1cm}
      
      {\Large \itshape 以时域为基，解析信号变化规律，\\以频域为核，挖掘信息本质特征。}\\[0.5cm]
      \rule{\textwidth}{1pt}\\[0.5cm]
      
      {\Huge \bfseries 信号分析与处理（AUTO2005）}\\[8pt]
      {\Huge \bfseries 课程笔记}\\[12pt]
      
      \rule{0.6\textwidth}{0.5pt}\\[1cm]
      
      {\Large \textbf{内容概要：} 信号描述 | 时域分析 | 频域分析 | 离散变换 | 滤波器设计}\\[8pt]
      {\Large \textbf{适用场景：} 课程复习 · 公式速查 }\\[1.5cm]
      
      {\Large \textbf{整理者：}duckweed}\\[6pt]
      {\large \textbf{GitHub：}https://github.com/Duckweed-yhb}\\[6pt]
      {\Large 整理日期：\today}\\[1cm]
      
      \rule{0.6\textwidth}{0.5pt}\\[1cm]
      
      {\large \textbf{版本：} v0.7 不完整版}\\[6pt]
      {\small \itshape Copyright © \the\year\ duckweed · 禁止商用 · 欢迎交流}\\[0.5cm]
      
      \rule{\textwidth}{1pt}
      \vfill
      \restoregeometry
    \end{titlepage}
}
% 生成封面
\makeSignalsCover

% 目录
\newpage
\thispagestyle{empty}      % 目录页无页码
\renewcommand{\contentsname}{目\quad 录}
\tableofcontents
\clearpage
\setcounter{page}{1}       % 正文从第1页开始编号

% 正文章节（按现有文件顺序引入）
\chapter{绪论}

\section{第一节：什么是信号}
信号是携带信息的物理量，是信息的载体，如电信号、光信号、声信号等，广泛存在于通信、控制、检测等工程领域。

\section{第二节：信号的描述与分类}
\subsection{一、信号的数学描述}
\begin{enumerate}
\item 数学表达式描述
  \begin{itemize}
  \item 核心：用函数关系 $f(t)$ / $f(n)$ 表征信号随时间/离散变量的变化规律
  \item 示例：正弦信号 $f(t) = A\sin(\omega t + \varphi)$、单位阶跃信号 $u(t)$
  \end{itemize}
\item 形态描述
  \begin{itemize}
  \item 核心：通过波形图、时域曲线等可视化方式描述信号特征
  \item 关键维度：幅值、周期、上升/下降沿、突变点等
  \end{itemize}
\end{enumerate}

\subsection{二、信号的分类}
\begin{enumerate}
\item 按信号中独立变量的个数
  \begin{itemize}
  \item 一维信号：仅含一个独立变量（如时间 $t$），如语音信号 $f(t)$
  \item 多维信号：含多个独立变量，如图像信号 $f(x,y)$、视频信号 $f(x,y,t)$
  \end{itemize}
\item 按信号的时间特性
  \begin{itemize}
  \item 确定信号：取值可由确定函数描述（如正弦信号、方波信号）
  \item 随机信号：取值无确定规律，仅能统计描述（如噪声、地震波）
  \end{itemize}
\item 按时间函数特性
  \begin{itemize}
  \item 连续时间信号：在连续时间范围内有定义（$t \in \mathbb{R}$），记为 $f(t)$
    \begin{itemize}
    \item 数学描述：$f(t), t \in (-\infty, +\infty)$，幅值可连续/离散（模拟信号是幅值也连续的连续时间信号）
    \end{itemize}
  \item 离散时间信号：仅在离散时间点有定义（$t = nT, n \in \mathbb{Z}$），记为 $f(n)$
    \begin{itemize}
    \item 数学描述：$f(n), n = 0, \pm1, \pm2, \dots$，也称为序列
    \end{itemize}
  \end{itemize}
\item 按周期性
  \begin{itemize}
  \item 周期信号：存在非零常数 $T$，满足 $f(t+T) = f(t)$（连续）/ $f(n+N) = f(n)$（离散）
    \begin{itemize}
    \item 判定规则：多个周期信号叠加，频率比为有理数时仍为周期信号
    \end{itemize}
  \item 非周期信号：无固定周期，如脉冲信号、衰减指数信号
  \end{itemize}
\item 按能量特点
  \subsubsection{信号的能量与功率分类}
  \begin{quote}
  以下定义基于 \textbf{1Ω 电阻} 负载，瞬时功率 $P = |f(t)|^2$
  \end{quote}
  
  \paragraph{1. 信号能量定义}
  \[
  E = \frac{1}{R} \int_{-\infty}^{\infty} |f(t)|^2 dt = \lim_{T \to \infty} \int_{-T}^{T} |f(t)|^2 dt
  \]
  \begin{quote}
  【说明】：$E$ 为信号总能量，单位为焦耳（J）
  \end{quote}
  
  \paragraph{2. 信号平均功率定义}
  \[
  P = \lim_{T \to \infty} \frac{1}{2T} \int_{-T}^{T} |f(t)|^2 dt = \lim_{T \to \infty} \frac{1}{T} \int_{-\frac{T}{2}}^{\frac{T}{2}} |f(t)|^2 dt
  \]
  \begin{quote}
  【说明】：$P$ 为信号平均功率，单位为瓦特（W）
  \end{quote}
  
  \paragraph{信号分类规则}
  根据信号的能量 $E$ 和功率 $P$ 取值，信号分为三类：
  \begin{itemize}
  \item \textbf{能量信号}
    \begin{quote}
    判定条件：$E < \infty$（能量有限），且 $P = 0$\\
    示例：矩形脉冲、衰减指数信号等\\
    【特点】：信号能量集中，随时间快速衰减，总能量有限
    \end{quote}
  \item \textbf{功率信号}
    \begin{quote}
    判定条件：$P < \infty$（功率有限），且 $E = \infty$\\
    示例：正弦信号、方波信号等周期信号\\
    【特点】：持续存在，功率稳定，但总能量无限
    \end{quote}
  \item \textbf{非能量非功率信号}
    \begin{quote}
    判定条件：$E = \infty$ 且 $P = \infty$\\
    示例：随时间无界增长的信号（如 $f(t) = t$）\\
    【特点】：既无有限总能量，也无稳定平均功率
    \end{quote}
  \end{itemize}
\end{enumerate}

\section{第三节：信号的分析与处理}
核心分析方法包括时域分析、频域分析、时频分析；处理手段包括滤波、降噪、特征提取、调制/解调等。

\section{第四节：自动控制系统中的信号分析与处理}
核心应用场景：
\begin{enumerate}
\item 控制系统中的信号类型：输入信号、反馈信号、误差信号
\item 分析方法：傅里叶变换分析系统频率响应、Z变换分析离散控制系统
\item 处理应用：噪声滤除、信号采样与重构、故障信号识别等
\end{enumerate}
% ========== 第一章 连续信号的分析 ==========
% 依据：26春_第1章_讲义-0415更新.pdf（223 页）
% 结构：1.1 时域描述和分析 / 1.2 频域分析 / 1.3 复频域分析

\chapter{连续信号的分析}

本章按"时域 $\to$ 频域 $\to$ 复频域"展开：先用基本函数与运算描述信号，
再把信号分解为正弦（复指数）分量得到频谱，最后引入衰减因子把频域推广为复频域。

\section{连续信号的时域描述和分析}

\subsection{连续信号的时域描述}

时域描述就是用\textbf{一个时间函数或一条曲线}表示信号随时间变化的特性。
基本函数（正弦、指数、阶跃、冲激、斜坡等）可以组合出许多更复杂的信号。

\subsubsection{正弦信号}

\begin{equation}
f(t)=K\sin(\omega t+\theta)=K\cos\Big(\omega t+\theta-\frac{\pi}{2}\Big),\qquad -\infty<t<\infty
\end{equation}
$K$ 为振幅，$\omega$ 为角频率，$\theta$ 为初始相角；正弦信号是周期信号，周期 $T=\dfrac{2\pi}{\omega}$。
余弦与正弦只在相位上相差 $\pi/2$，通常也归为正弦信号。性质：

\begin{enumerate}
\item 两个\textbf{同频率}的正弦信号相加，结果仍是原频率的正弦信号；
\item 若 $f_1=nf_0$（$n$ 为整数），则合成信号是频率为 $f_0$ 的\textbf{非正弦}周期信号：
  $f_0$ 为\textbf{基波频率}，$f_1=nf_0$ 为 $n$ 次\textbf{谐波频率}，周期信号可分解为基波与一系列谐波；
\item 正弦信号的微分和积分仍是同频率的正弦信号；
\item 两个正弦（周期）信号相加，若周期之比 $T_1/T_2$ 为\textbf{有理数}，则仍是周期信号，
  周期为 $T_1$、$T_2$ 的\textbf{最小公倍数}。
\end{enumerate}

第 4 条即判断方法：读出各分量的周期，求周期之比；有理数则为周期信号（周期取最小公倍数），
无理数则合成信号不是周期信号。

\subsubsection{指数信号}

$f(t)=K\mathrm{e}^{at}$：$a<0$ 为衰减信号，$a=0$ 为直流信号，$a>0$ 为增长信号；$|a|$ 反映变化快慢。

\subsubsection{复指数信号}

\begin{equation}
f(t)=K\mathrm{e}^{st}=K\mathrm{e}^{\sigma t}\big(\cos\omega t+\mathrm{j}\sin\omega t\big)
=\operatorname{Re}[f(t)]+\mathrm{j}\operatorname{Im}[f(t)],\qquad s=\sigma+\mathrm{j}\omega,
\end{equation}
其中 $\operatorname{Re}[f(t)]=K\mathrm{e}^{\sigma t}\cos\omega t$，
$\operatorname{Im}[f(t)]=K\mathrm{e}^{\sigma t}\sin\omega t$。按 $\sigma$、$\omega$ 讨论：

\begin{center}
\begin{tabular}{lll}
\toprule
参数 & 信号形态 & 说明 \\
\midrule
$\sigma=0,\ \omega=0$ & 直流信号 & 常值 \\
$\sigma=0,\ \omega\neq0$ & 正弦信号 & 实部余弦、虚部正弦 \\
$\sigma\neq0,\ \omega=0$ & （实）指数信号 & $\sigma<0$ 收敛，$\sigma>0$ 发散 \\
$\sigma\neq0,\ \omega\neq0$ & 指数包络的正弦型振荡 & $\sigma<0$ 收敛，$\sigma>0$ 发散 \\
\bottomrule
\end{tabular}
\end{center}

实际信号总是实的，复指数信号不可能实际产生；其意义在于\textbf{实部与虚部表示了指数包络的正弦型振荡}，
从而把直流、指数、正弦型以及带包络的正弦型信号统一表示：
\begin{equation}
A\cos(\omega t+\varphi)=A\operatorname{Re}\big[\mathrm{e}^{\mathrm{j}(\omega t+\varphi)}\big],\qquad
A\sin(\omega t+\varphi)=A\operatorname{Im}\big[\mathrm{e}^{\mathrm{j}(\omega t+\varphi)}\big].
\end{equation}

\subsubsection{奇异信号}

用奇异函数表示的连续时间信号，\textbf{函数本身或其导数（含高阶导数）有不连续点}，
是从实际信号中抽象出来的典型信号。

\paragraph{（1）单位斜坡信号}
\begin{equation}
r(t)=
\begin{cases}
t, & t\ge0,\\
0, & t<0,
\end{cases}
\qquad
r(t-t_0)=
\begin{cases}
t-t_0, & t\ge t_0,\\
0, & t<t_0,
\end{cases}
\qquad r'(t)=u(t),
\end{equation}
其导数在 $t=0$ 处不连续。用斜坡信号可表示"截平的"斜变信号与三角形脉冲（此时要增加对斜率的计算）。

\paragraph{（2）单位阶跃信号}
\begin{equation}
u(t)=
\begin{cases}
1, & t>0,\\
0, & t<0,
\end{cases}
\qquad
u(t-t_0)=
\begin{cases}
1, & t>t_0,\\
0, & t<t_0,
\end{cases}
\qquad u(0)=\tfrac12,
\end{equation}
可用来描述信号的\textbf{接入特性}（如接入单位直流电源）。矩形脉冲
\begin{equation}
G_T(t)=u\Big(t+\frac{T}{2}\Big)-u\Big(t-\frac{T}{2}\Big),
\qquad
\frac{\mathrm{d}r(t)}{\mathrm{d}t}=u(t).
\end{equation}

\paragraph{（3）单位冲激信号（Dirac $\delta$ 函数）}
持续时间无穷小、瞬间幅度无穷大、涵盖面积恒为 $1$ 的理想函数（分布）：
\begin{equation}
\delta(t)=0\ (t\neq0),\qquad
\int_{-\infty}^{\infty}\delta(t)\,\mathrm{d}t=\int_{0^-}^{0^+}\delta(t)\,\mathrm{d}t=1,
\end{equation}
\begin{equation}
\delta(t)=\lim_{\tau\to0}\frac{1}{\tau}\Big[u\Big(t+\frac{\tau}{2}\Big)-u\Big(t-\frac{\tau}{2}\Big)\Big]
=\lim_{\tau\to0}\frac{1}{\tau}\Big(1-\frac{|t|}{\tau}\Big)\big[u(t+\tau)-u(t-\tau)\big].
\end{equation}
物理上表示理想化的质点或点电荷的密度。性质：
\begin{itemize}
\item \textbf{抽样（筛选）特性}：$f(t)$ 在 $t=0$ 连续时
  $\displaystyle\int_{-\infty}^{\infty}f(t)\delta(t)\mathrm{d}t=f(0)$；在 $t=t_0$ 连续时
  $\displaystyle\int_{-\infty}^{\infty}f(t)\delta(t-t_0)\mathrm{d}t=f(t_0)$；
\item \textbf{偶函数特性}：$\delta(t)=\delta(-t)$；
\item \textbf{与阶跃互为积分和微分}：
  $\displaystyle\int_{-\infty}^{t}\delta(\tau)\mathrm{d}\tau=u(t)$，
  $\dfrac{\mathrm{d}u(t)}{\mathrm{d}t}=\delta(t)$。
\end{itemize}
注意：$\delta$ 函数只有出现在积分以内才有实质意义。

\paragraph{（4）单位冲激偶信号}
$\delta'(t)$ 是冲激函数的微分。$\tau$ 可从负值或正值趋于 $0$，故微分后得到正、负一对冲激。性质：
\begin{itemize}
\item $\displaystyle\int_{-\infty}^{\infty}\delta'(t)f(t)\mathrm{d}t=-f'(0)$（$f'$ 在 $0$ 点连续）；
\item $\displaystyle\int_{-\infty}^{\infty}\delta'(t-t_0)f(t)\mathrm{d}t=-f'(t_0)$（$f'$ 在 $t_0$ 连续）；
\item 面积为零：$\displaystyle\int_{-\infty}^{\infty}\delta'(t)\mathrm{d}t=0$；
\item 奇函数：$\delta'(-t)=-\delta'(t)$。
\end{itemize}

\subsection{连续信号的时域运算}

四类运算：基本运算（尺度变换、翻转、平移）、叠加和相乘、微分和积分、卷积运算及其性质。
卷积运算要求熟练掌握。

\subsubsection{基本运算}

\begin{itemize}
\item \textbf{幅度尺寸变换} $f_2(t)=af_1(t)$：对幅度放缩，\textbf{不改变}信号基本特征；
\item \textbf{时间尺寸变换} $f_2(t)=f_1(at)$：\textbf{改变}信号基本特征，频谱随之变化；
\item \textbf{翻转} $f(-t)$：以纵坐标轴为中心对称映射；
\item \textbf{平移}：$t_0>0$ 时 $f(t)\to f(t-t_0)$ 为\textbf{滞后}（右移），$f(t)\to f(t+t_0)$ 为\textbf{超前}（左移）。
\end{itemize}
复合运算（如求 $f(at\pm t_0)$）按"先尺度/翻转、后平移"逐步代入，或先整理成 $f\big(a(t\pm t_1)\big)$ 再作图。

\subsubsection{叠加与相乘}

$f(t)=f_1(t)+f_2(t)$ 为相同时刻的代数和；$f(t)=f_1(t)f_2(t)$ 为相同时刻的乘积（常用于调制、加窗截取）。

\subsubsection{微分与积分}

微分 $f'(t)$ \textbf{突出}信号的变化部分；积分 $\displaystyle\int_{-\infty}^{t}f(\tau)\mathrm{d}\tau$
\textbf{平滑}信号的突变部分。奇异信号间可串成一条链：
\begin{equation}
r(t)\xrightarrow{\ \text{微分}\ }u(t)\xrightarrow{\ \text{微分}\ }\delta(t)
\xrightarrow{\ \text{微分}\ }\delta'(t),\qquad
\int_{-\infty}^{t}u\mathrm{d}\tau=r(t),\quad
\int_{-\infty}^{t}\delta\mathrm{d}\tau=u(t),\quad
\int_{-\infty}^{t}\delta'\mathrm{d}\tau=\delta(t).
\end{equation}

\subsubsection{卷积运算}

\begin{equation}
x_1(t)*x_2(t)=\int_{-\infty}^{\infty}x_1(\tau)x_2(t-\tau)\,\mathrm{d}\tau
=\int_{-\infty}^{\infty}x_2(\tau)x_1(t-\tau)\,\mathrm{d}\tau .
\end{equation}
\textbf{图解法步骤}：变量替换（$t\to\tau$）$\to$ 翻转（$x_2(\tau)\to x_2(-\tau)$）$\to$
平移（$\to x_2(t-\tau)$）$\to$ 相乘 $\to$ 积分。
关键在\textbf{积分上下限的选取}：它取决于两个函数图形交叠部分的范围，需按 $t$ 分段讨论
（不交叠时为零，部分交叠时按交叠端点积分，完全包含时按非零区间积分）。

\textbf{性质}：
\begin{itemize}
\item \textbf{交换律}：$x_1*x_2=x_2*x_1$（令 $\lambda=t-\tau$ 即证）；
\item \textbf{分配律}：$x*(x_1+x_2)=x*x_1+x*x_2$；
\item \textbf{结合律}：$(x*x_1)*x_2=x*(x_1*x_2)$（改变积分次序并换元）；
\item \textbf{卷积的微分}：$\dfrac{\mathrm{d}}{\mathrm{d}t}(x_1*x_2)=\dfrac{\mathrm{d}x_1}{\mathrm{d}t}*x_2
  =x_1*\dfrac{\mathrm{d}x_2}{\mathrm{d}t}$——与乘法的微分不同，\textbf{只需对其中一个函数微分}；
\item \textbf{卷积的积分}：$\displaystyle\int_{-\infty}^{t}\big[x_1(\lambda)*x_2(\lambda)\big]\mathrm{d}\lambda
  =x_1(t)*\int_{-\infty}^{t}x_2\mathrm{d}\lambda=\int_{-\infty}^{t}x_1\mathrm{d}\lambda*x_2(t)$；
\item \textbf{与冲激、冲激偶、阶跃的卷积}：
  \begin{equation}
  \begin{gathered}
  x(t)*\delta(t)=x(t),\qquad x(t)*\delta(t-t_0)=x(t-t_0),\qquad
  x(t-t_1)*\delta(t-t_2)=x(t-t_1-t_2),\\
  \delta(t-t_1)*\delta(t-t_2)=\delta(t-t_1-t_2),\qquad
  x(t)*\delta'(t)=x'(t),\qquad
  x(t)*u(t)=\int_{-\infty}^{t}x(\tau)\mathrm{d}\tau,
  \end{gathered}
  \end{equation}
  推广：$x(t)*\delta^{(k)}(t)=x^{(k)}(t)$，$x(t)*\delta^{(k)}(t-t_0)=x^{(k)}(t-t_0)$
  （$k$ 为正整数表示导数阶次，为负整数表示重积分次数）。
\end{itemize}

\textbf{利用性质简化运算}：把一个信号微分、另一个积分，可把卷积化为冲激与斜坡等简单信号的卷积，
避免分段积分。卷积结果的\textbf{时宽}等于两个函数各自时宽之和。

\textbf{理解}：卷积实质上是对信号进行\textbf{滤波}；物理上可代表某系统对物理量或输入的调制或污染；
应用中可把非光滑函数光滑化；离散卷积在数字图像处理中常用。
注：现实中信号通常只在 $t\ge0$ 时有意义，此时亦需 $t-\tau\ge0$，故 $\tau\in[0,t]$，
自动控制方向的教材常写成 $x_1(t)*x_2(t)=\displaystyle\int_{0}^{t}x_1(\tau)x_2(t-\tau)\mathrm{d}\tau$。

\subsection{连续信号的时域分解}

思路：把复杂信号分解成一些简单信号或基本信号，再由简单信号的结果叠加。

\subsubsection{直流分量与交流分量}

$f(t)=f_D+f_A(t)$，$f_D$ 为直流分量（平均值），$f_A(t)$ 为交流分量。

\subsubsection{偶分量与奇分量}

\begin{equation}
f_e(t)=\tfrac12\big[f(t)+f(-t)\big]=f_e(-t),\qquad
f_o(t)=\tfrac12\big[f(t)-f(-t)\big]=-f_o(-t).
\end{equation}
任何信号都可分解为偶分量与奇分量之和（由 $f(-t)=f_e(t)-f_o(t)$ 可验证）。

\subsubsection{分解成冲激函数之和}

用窄矩形脉冲近似 $f(t)$，令 $\Delta t\to0$：
\begin{equation}
\lim_{\Delta t\to0}\frac{1}{\Delta t}\big\{u(t-k\Delta t)-u[t-(k+1)\Delta t]\big\}=\delta(t-\tau)
\ \Longrightarrow\
f(t)=\int_{-\infty}^{\infty}f(\tau)\delta(t-\tau)\mathrm{d}\tau=f(t)*\delta(t).
\end{equation}
即任意信号都可分解为一系列具有不同强度的冲激函数。

\subsubsection{从向量空间到函数空间}

二维中向量 $x$ 在 $y$ 上的投影系数 $c=\dfrac{\langle x,y\rangle}{\langle y,y\rangle}$
使 $\min\|x-cy\|_2^2$ 成立（误差函数为凸函数，导数为零处最小）；
$n$ 维中若 $V_1,\dots,V_n$ 正交，则 $x=\sum c_iV_i$ 的系数唯一：
\begin{equation}
c_i=\frac{\langle x,V_i\rangle}{\langle V_i,V_i\rangle},\qquad i=1,2,\dots,n,
\end{equation}
且 $n$ 维向量必须用 $n$ 个正交分量表示，否则产生误差。
把维数推广到无穷并限制向量长度有限，得到\textbf{希尔伯特（Hilbert）空间}——
完备的内积空间，其元素（向量）在实际应用中可代表一列复数或一个函数，于是函数也可在其中计算内积。

\subsubsection{正交函数集与完备正交函数集}

两个非零函数在 $(t_1,t_2)$ 内\textbf{正交}指（上标 $*$ 为共轭，实函数时 $f_2^*=f_2$）
\begin{equation}
\int_{t_1}^{t_2}f_1(t)f_2^{*}(t)\,\mathrm{d}t=0;
\end{equation}
函数集 $\{f_i(t)\}$ 为\textbf{正交函数集}指
\begin{equation}
\int_{t_1}^{t_2}f_i(t)f_j^{*}(t)\,\mathrm{d}t=
\begin{cases}
0, & i\neq j,\\
k_i, & i=j,
\end{cases}
\end{equation}
若进一步不存在非零函数 $\varphi(t)$ 使
$\displaystyle\int_{t_1}^{t_2}f_i(t)\varphi^{*}(t)\mathrm{d}t=0\ (i=1,\dots,n)$，则为\textbf{完备正交函数集}。
常见的完备正交函数集（在 $(t_0,t_0+T)$ 上，$T=2\pi/\omega_0$）：
\begin{itemize}
\item 三角函数集 $\big\{1,\cos\omega_0t,\cos2\omega_0t,\dots,\sin\omega_0t,\sin2\omega_0t,\dots\big\}$；
\item 复指数函数集 $\big\{\mathrm{e}^{\mathrm{j}n\omega_0t}\big\}$，$n=0,\pm1,\pm2,\dots$；
\item 其它：勒让德函数集、切比雪夫多项式集合、沃尔什函数集等（了解即可）。
\end{itemize}

\subsubsection{信号的正交分解}

用 $n$ 个正交函数近似 $x(t)$：$x(t)=\sum_{i=1}^{n}c_if_i(t)+x_e(t)$，$t\in(t_1,t_2)$，
使\textbf{均方误差}最小：
\begin{equation}
\overline{x_e^2(t)}=\frac{1}{t_2-t_1}\int_{t_1}^{t_2}
\Big[x(t)-\sum_{i=1}^{n}c_if_i(t)\Big]^2\mathrm{d}t,
\end{equation}
最优系数（即信号在正交分量上的\textbf{投影}）
\begin{equation}
c_i=\frac{\displaystyle\int_{t_1}^{t_2}x(t)f_i(t)\,\mathrm{d}t}
{\displaystyle\int_{t_1}^{t_2}\big[f_i(t)\big]^2\mathrm{d}t}
=\frac{1}{k_i}\int_{t_1}^{t_2}x(t)f_i(t)\,\mathrm{d}t,
\qquad
\text{复变函数情形 }c_i=\frac{\int x(t)f_i^{*}(t)\mathrm{d}t}{\int\big|f_i(t)\big|^2\mathrm{d}t}.
\end{equation}
两种推导思路：对 $\overline{x_e^2(t)}$ 求 $c_i$ 的偏导并令其为零（优化角度）；
或直接取 $x_e(t)=0$ 的最优情况，两边与 $f_i(t)$ 取内积，由正交性立即解出 $c_i$（内积/投影角度）。
$c_i$ 只与 $x(t)$ 和 $f_i(t)$ 有关，与其余分量无关。

\subsubsection{帕斯瓦尔方程}

若 $\overline{x_e^2(t)}=0$，则 $x(t)$ 可由 $n$ 个正交函数精确描述，正交函数集完备。
一般 $\overline{x_e^2(t)}>0$，$n$ 越大均方误差越小；$n\to\infty$ 时 $\overline{x_e^2(t)}=0$，且
\begin{equation}
\int_{t_1}^{t_2}x^2(t)\,\mathrm{d}t=\sum_{i=1}^{\infty}c_i^2k_i .
\end{equation}
\textbf{帕斯瓦尔（Parseval）方程}表示信号分解的能量关系：信号 $x(t)$ 的能量等于此信号在完备正交函数集中
各分量能量之和。

\section{连续信号的频域分析}

\textbf{频域分析法}：三角函数集是完备正交函数集，任意信号都可分解为三角函数表达式；
任何信号都可视为一系列正弦信号的组合，其频率、相位反映原信号的性质。
用频率域的特性（频率特性）描述时间域信号，具有很强的物理意义——频率特性是信号的客观性质
（光的颜色由频率决定，音调的不同也在于频率差异）。

傅里叶变换起源于传热学问题（铁环上温度的不规则分布可用许多正弦分量的频率描述）。
发展脉络：1928 Nyquist（抽样定理）、1935 Wiener（功率谱）、1948 Shannon（信息论）、
1965 Cooley 与 Tukey（FFT）、1984 Morlet 等（小波变换）。

\subsection{周期信号的傅里叶级数}

\subsubsection{狄里赫利条件}

周期信号能展开为傅里叶级数的\textbf{充分条件}：一个周期内间断点有限个；极大值、极小值有限个；
信号绝对可积，即 $\displaystyle\int_{t_0}^{t_0+T_1}|f(t)|\mathrm{d}t$ 为有限值。
通常实际的周期信号都满足此条件。

\subsubsection{三角函数形式的傅里叶级数}

设周期为 $T_1$、基波角频率 $\omega_1=2\pi/T_1$，
\begin{equation}
f(t)=\frac{a_0}{2}+\sum_{n=1}^{\infty}\big(a_n\cos n\omega_1t+b_n\sin n\omega_1t\big),
\end{equation}
\begin{align}
a_0&=\frac{2}{T_1}\int_{-\frac{T_1}{2}}^{\frac{T_1}{2}}f(t)\mathrm{d}t,\\
a_n&=\frac{2}{T_1}\int_{-\frac{T_1}{2}}^{\frac{T_1}{2}}f(t)\cos n\omega_1t\,\mathrm{d}t\quad(n\ \text{的偶函数}),\\
b_n&=\frac{2}{T_1}\int_{-\frac{T_1}{2}}^{\frac{T_1}{2}}f(t)\sin n\omega_1t\,\mathrm{d}t\quad(n\ \text{的奇函数}).
\end{align}
$a_0/2$ 为直流分量，$a_n$、$b_n$ 分别为余弦、正弦分量的幅度。若不整理成首项 $a_0/2$ 的形式，也可写成
$f(t)=a_0+\sum_{n=1}^{\infty}(\cdot)$，对应 $a_0=\frac{1}{T_1}\int f\mathrm{d}t$，两者本质相同。

\textbf{系数求解}两法：① 正交分解法——用内积求投影，分母
$\int\cos^2n\omega_1t\,\mathrm{d}t=T_1/2$（用半角公式 $\cos^2\alpha=\frac{1+\cos2\alpha}{2}$、
$\sin^2\alpha=\frac{1-\cos2\alpha}{2}$）；② 积分法——展开式两边同乘 $\sin k\omega_1t$ 后在一个周期内积分，
由正交性只剩 $n=k$ 项，得 $b_n=\frac{2}{T_1}\int f(t)\sin n\omega_1t\,\mathrm{d}t$。

\subsubsection{余弦形式}

\begin{equation}
f(t)=\frac{c_0}{2}+\sum_{n=1}^{\infty}c_n\cos(n\omega_1t+\varphi_n),\qquad
c_0=a_0,\quad c_n=\sqrt{a_n^2+b_n^2},\quad \varphi_n=-\arctan\frac{b_n}{a_n},
\end{equation}
且 $a_n=c_n\cos\varphi_n$、$b_n=-c_n\sin\varphi_n$；$c_n$ 是 $n$ 的偶函数，$\varphi_n$ 是 $n$ 的奇函数。
即周期信号可分解为直流分量和一系列余弦形式的交流分量，直流分量
$\dfrac{a_0}{2}=\dfrac{1}{T_1}\displaystyle\int_{-T_1/2}^{T_1/2}f(t)\mathrm{d}t$ 是信号在一个周期内的平均值。

\subsubsection{周期信号的频谱}

$n=1$ 的分量为\textbf{基波分量}（一次谐波），$n=2,3,\dots$ 为二次、三次……谐波分量。
$c_n$ 随频率的分布为\textbf{幅度频谱（幅度谱）}，$\varphi_n$ 随频率的分布为\textbf{相位频谱（相位谱）}，
谱线顶点可用\textbf{包络线}连接。谱线只出现在 $0,\omega_1,2\omega_1,\dots$ 等\textbf{离散}频率点上，
故为\textbf{离散谱}，谱线间隔 $\omega_1=2\pi/T_1$。

\textbf{如何理解相位}：海浪可用不同波长正弦波的线性组合近似；三种波列叠加时，
是否形成图中虚线所示的巨大波浪，取决于这些正弦分量的\textbf{相对相位}。

\subsubsection{复指数（指数）形式的傅里叶级数}

把三角形式用欧拉公式 $\mathrm{e}^{\mathrm{j}n\omega_1t}=\cos n\omega_1t+\mathrm{j}\sin n\omega_1t$ 逐项改写
（$a_n$ 偶、$b_n$ 奇），得
\begin{equation}
f(t)=\sum_{n=-\infty}^{\infty}F_n\mathrm{e}^{\mathrm{j}n\omega_1t},\qquad
F_n=\frac{a_n-\mathrm{j}b_n}{2}=\frac{1}{T_1}\int_{-\frac{T_1}{2}}^{\frac{T_1}{2}}f(t)\mathrm{e}^{-\mathrm{j}n\omega_1t}\mathrm{d}t,
\end{equation}
\begin{equation}
F_n=|F_n|\mathrm{e}^{\mathrm{j}\varphi_n}=\frac{c_n}{2}\mathrm{e}^{\mathrm{j}\varphi_n},\qquad
F_{-n}=\frac{a_n+\mathrm{j}b_n}{2}=F_n^{*},
\end{equation}
\begin{equation}
F_n+F_{-n}=a_n,\qquad |F_n|=\frac{c_n}{2},\qquad b_n=\mathrm{j}(F_n-F_{-n}).
\end{equation}
即指数形式的\textbf{幅度变成三角形式的一半}，正负频率的系数\textbf{互为共轭}。
\textbf{负频率}是数学运算的结果，把负频率项与相应正频率项成对合并才是实际的频谱函数
（了解：正角频率对应向量逆时针旋转、负的顺时针；多普勒频移的正负表示接近或远离；
电机中正负角频率表示转向；通信中基于希尔伯特变换抵消负频率以节省一半带宽）。

\subsubsection{单边谱与双边谱}

单边谱（$n=0,1,2,\dots$）即三角形式 $\dfrac{c_0}{2}+\sum c_n\cos(n\omega_1t+\varphi_n)$；
双边谱（$n=0,\pm1,\pm2,\dots$）即指数形式 $\sum F_n\mathrm{e}^{\mathrm{j}n\omega_1t}$，$|F_n|=c_n/2$。
$|F_n|$ 是 $n$ 的偶函数、$\varphi_n$ 是 $n$ 的奇函数；当 $F_n$ 为实数时，
可用 $F_n$ 的正、负表示 $\varphi_n$ 的 $0$ 或 $\pi$，从而把幅度与相位画在一张图上。

\subsubsection{周期信号的功率分配与功率频谱}

周期信号（功率信号）的平均功率为 $P=\dfrac{1}{T_1}\displaystyle\int_{-T_1/2}^{T_1/2}|f(t)|^2\mathrm{d}t$，
由帕斯瓦尔公式
\begin{equation}
P=\frac{1}{T_1}\int_{-\frac{T_1}{2}}^{\frac{T_1}{2}}\big|f(t)\big|^2\mathrm{d}t
=\frac{c_0^2}{4}+\frac12\sum_{n=1}^{\infty}c_n^2=\sum_{n=-\infty}^{\infty}\big|F_n\big|^2,
\end{equation}
即周期信号在时域的平均功率等于直流、基波及各次谐波平均功率之和（时域与频域的能量守恒）。
把各次谐波（含直流）的平均功率分配表示成谱线形式，即\textbf{功率频谱}。

\subsubsection{函数的对称性与傅里叶系数的关系}

\textbf{偶函数} $f(t)=f(-t)$：
\begin{equation}
a_0=\frac{4}{T_1}\int_{0}^{\frac{T_1}{2}}f(t)\mathrm{d}t,\qquad
a_n=\frac{4}{T_1}\int_{0}^{\frac{T_1}{2}}f(t)\cos n\omega_1t\,\mathrm{d}t,\qquad b_n=0,
\end{equation}
\begin{equation}
c_n=|a_n|=2|F_n|,\qquad F_n=F_{-n}=\frac{a_n}{2}\ (\text{实部，偶函数}),\qquad \varphi_n=0\ \text{或}\ \pi
\end{equation}
（由 $F_n=\frac{a_n}{2}=|F_n|\mathrm{e}^{\mathrm{j}\varphi_n}$ 得 $\mathrm{e}^{\mathrm{j}\varphi_n}=\pm1$，
故 $\cos\varphi_n=\pm1$、$\sin\varphi_n=0$，取 $\varphi_n\in(-\pi,\pi]$ 即 $0$ 或 $\pi$）。

\textbf{奇函数} $f(t)=-f(-t)$：
\begin{equation}
a_0=0,\quad a_n=0,\quad b_n=\frac{4}{T_1}\int_{0}^{\frac{T_1}{2}}f(t)\sin n\omega_1t\,\mathrm{d}t,\quad
c_n=|b_n|,\quad F_n=-F_{-n}=-\frac{\mathrm{j}b_n}{2},\quad \varphi_n=\mp\frac{\pi}{2}.
\end{equation}
注意：奇函数叠加直流分量后 $a_0\neq0$、整体不再是奇函数，但余弦项系数 $a_n$ 仍为零
（常数在一个完整周期上与 $\cos n\omega_1t$ 正交）。

\textbf{奇谐函数（半波对称）} $f(t)=-f\big(t\pm\frac{T_1}{2}\big)$（波形平移半周期并上下反转）：
\begin{equation}
a_0=0,\qquad a_n=b_n=0\ (n\ \text{为偶数}),\qquad
a_n,b_n=\frac{4}{T_1}\int_{0}^{\frac{T_1}{2}}f(t)\begin{cases}\cos\\ \sin\end{cases}n\omega_1t\,\mathrm{d}t
\ (n\ \text{为奇数}),
\end{equation}
即只含\textbf{奇次}谐波的正弦、余弦项。

\subsubsection{周期信号的傅里叶近似与吉布斯现象}

对展开式截断，取前 $N$ 项逼近（以方波为例 $f_N(t)=\sum_{n=1}^{N}b_{2n-1}\sin(2n-1)\omega_1t$）：
\begin{equation}
\varepsilon_N(t)=f(t)-f_N(t),\qquad
E_N=\overline{\varepsilon_N^2(t)}=\frac{1}{T_1}\int_{t_0}^{t_0+T_1}\varepsilon_N^2(t)\,\mathrm{d}t .
\end{equation}
\begin{itemize}
\item 项数越多，叠加波形越逼近原信号，\textbf{均方误差越小}；$N\to\infty$ 时 $f_N(t)\to f(t)$；
\item \textbf{低频分量影响脉冲顶部}（低频越多，信号变化越缓慢）；
  \textbf{高频分量影响脉冲跳变沿}（高频越多，信号变化越剧烈）；
\item \textbf{吉布斯现象}：项数越多，合成波形中的峰起越靠近 $f(t)$ 的不连续点；$N$ 很大时峰起值趋于常数，
  约等于总跳变值的 $9\%$，并从不连续点开始以起伏振荡形式衰减（1848 年 Wilbraham 提出、1899 年 Gibbs 证明；
  力学中也称"振铃现象"）。
\end{itemize}

\subsubsection{典型周期信号的傅里叶级数}

\paragraph{（1）周期矩形脉冲信号}
脉宽 $\tau$、幅度 $E$、周期 $T_1$：
\begin{equation}
f(t)=E\Big[u\Big(t+\frac{\tau}{2}\Big)-u\Big(t-\frac{\tau}{2}\Big)\Big],\quad
t\in\Big[-\frac{T_1}{2},\frac{T_1}{2}\Big],\ \tau\in(0,T_1),
\end{equation}
它是偶函数，故 $b_n=0$、$c_n=a_n$：
\begin{equation}
a_0=\frac{2E\tau}{T_1},\qquad
a_n=\frac{2E}{n\pi}\sin\frac{n\pi\tau}{T_1}=\frac{2E\tau}{T_1}\mathrm{Sa}\Big(\frac{n\omega_1\tau}{2}\Big),\qquad
F_n=\frac{a_n}{2}=\frac{E\tau}{T_1}\mathrm{Sa}\Big(\frac{n\omega_1\tau}{2}\Big),
\end{equation}
其中取样函数（Sampling function，又译抽样/采样函数）
$\mathrm{Sa}(x)=\dfrac{\sin x}{x}$，$\mathrm{Sa}(0)=1$，$\displaystyle\int_{-\infty}^{\infty}\mathrm{Sa}(x)\mathrm{d}x=\pi$。

\textbf{频谱特点}（$F_n=\frac{E\tau}{T_1}\mathrm{Sa}(\frac{n\omega_1\tau}{2})$）：
\begin{enumerate}
\item 谱线离散，间隔 $\omega_1=2\pi/T_1$；
\item 各分量大小正比于 $E$、$\tau$，反比于 $T_1$；直流分量峰值为 $E\tau/T_1$；
\item 各谱线幅度按 $\mathrm{Sa}\big(\frac{\omega\tau}{2}\big)$ 包络变化，当 $\omega=\frac{2m\pi}{\tau}$
  （$m=\pm1,\pm2,\dots$）时包络过零，\textbf{第一个过零点频率}为 $\frac{2\pi}{\tau}$；
\item 谱线整体上有减小趋势，较高幅值的谱线集中在第一个过零点 $\omega=\frac{2\pi}{\tau}$ 以内（主要能量）；
  矩形信号的\textbf{频带宽度} $B_\omega=\frac{2\pi}{\tau}$（或 $B_f=\frac{1}{\tau}$），只与 $\tau$ 成反比；
\item 允许一定失真时，可只传送 $\omega\le\frac{2\pi}{\tau}$ 内的分量，舍弃更高频率分量。
\end{enumerate}
谐波性：$T_1$ 不变、$\tau$ 变小则第一个过零点频率增大、谱线幅值减小；$\tau$ 不变、$T_1$ 增大则过零点频率不变、
振幅变小、谱线越来越密集；$T_1\to\infty$ 时成为非周期信号，谱线间隔 $\to0$（连续频谱）。

\paragraph{（2）对称方波信号}
取 $\tau=T_1/2$，令 $\bar f(t)=f(t)-\frac{E}{2}$（向下平移 $E/2$），则
$a_0=0$、$b_n=0$，
\begin{equation}
a_n=\frac{2E}{n\pi}\sin\frac{n\pi\tau}{T_1}=E\,\mathrm{Sa}\Big(\frac{n\pi}{2}\Big),\qquad
\bar f(t)=\frac{2E}{\pi}\sum_{n=1}^{\infty}\frac{1}{n}\sin\frac{n\pi}{2}\cos(n\omega_1t).
\end{equation}
它仍是偶函数，且是\textbf{奇谐函数}（只含奇次谐波）；其复傅里叶系数与周期矩形脉冲相同
（$\bar F_n=F_n$，因为常数 $E/2$ 在 $n\neq0$ 时与 $\mathrm{e}^{-\mathrm{j}n\omega_1t}$ 正交），
故对称方波是周期矩形脉冲的\textbf{一类特殊情况}。

\paragraph{（3）复指数信号与正余弦信号}
\begin{equation}
F(n\omega_1)=\frac{1}{T_1}\int_{-\frac{T_1}{2}}^{\frac{T_1}{2}}
\mathrm{e}^{\mathrm{j}\omega_1t}\mathrm{e}^{-\mathrm{j}n\omega_1t}\mathrm{d}t
=\mathrm{Sa}\big((1-n)\pi\big)=
\begin{cases}
1, & n=1,\\
0, & n\neq1,
\end{cases}
\end{equation}
即 $\mathrm{e}^{\mathrm{j}\omega_1t}$ 仅在 $\omega_1$ 处有幅度为 $1$ 的分量（$\varphi_n=0$），
说明复指数信号是正弦信号的一种表现形式。由
$\cos\omega_1t=\frac12(\mathrm{e}^{\mathrm{j}\omega_1t}+\mathrm{e}^{-\mathrm{j}\omega_1t})$、
$\sin\omega_1t=\frac{1}{2\mathrm{j}}(\mathrm{e}^{\mathrm{j}\omega_1t}-\mathrm{e}^{-\mathrm{j}\omega_1t})$ 得
二者在 $\pm\omega_1$ 处幅度均为 $\frac12$，但正弦的相位比余弦\textbf{滞后} $\pi/2$。

\subsection{非周期信号的傅里叶变换}

\subsubsection{从傅里叶级数到傅里叶变换}

非周期信号可看作\textbf{周期无穷大}的周期信号（基波频率无穷小）。$T_1$ 增大时周期 $\uparrow$、幅度 $\downarrow$、
谱线密集程度 $\uparrow$；$T_1\to\infty$ 时谱线连续，而谱线幅度变为\textbf{无穷小量}。
为避免无穷小量，把 $f(t)$ 周期延拓为 $\hat f(t)$ 展开，并定义频谱密度 $F(\omega)=T_1F(n\omega_1)$：
\begin{equation}
F(\omega)=\lim_{T_1\to\infty}\int_{-\frac{T_1}{2}}^{\frac{T_1}{2}}f(t)\mathrm{e}^{-\mathrm{j}n\omega_1t}\mathrm{d}t
=\int_{-\infty}^{\infty}f(t)\mathrm{e}^{-\mathrm{j}\omega t}\mathrm{d}t
\end{equation}
（$T_1\to\infty$ 时 $\omega_1=2\pi/T_1\to0$，$\omega_1\to\mathrm{d}\omega$、$n\omega_1\to\omega$）。
反变换由 $\hat f(t)=\sum\hat F_n\mathrm{e}^{\mathrm{j}n\omega_1t}$ 取极限（用 $\frac{1}{T_1}=\frac{\omega_1}{2\pi}$、
$\sum\to\int$）得
\begin{equation}
f(t)=\frac{1}{2\pi}\int_{-\infty}^{\infty}F(\omega)\mathrm{e}^{\mathrm{j}\omega t}\,\mathrm{d}\omega .
\end{equation}

\subsubsection{傅里叶变换对与频谱密度函数}

\begin{equation}
F(\omega)=\mathcal{F}\big[f(t)\big]=\int_{-\infty}^{\infty}f(t)\mathrm{e}^{-\mathrm{j}\omega t}\mathrm{d}t,\qquad
f(t)=\mathcal{F}^{-1}\big[F(\omega)\big]=\frac{1}{2\pi}\int_{-\infty}^{\infty}F(\omega)\mathrm{e}^{\mathrm{j}\omega t}\mathrm{d}\omega,
\end{equation}
记作 $f(t)\overset{\mathcal{F}}{\longleftrightarrow}F(\omega)$。$F(\omega)$ 称为\textbf{频谱密度函数}（简称频谱函数）：
频谱密度 $F(n\omega_1)/\omega_1$ 表示\textbf{单位频带}的频谱值，小矩形面积等于 $\omega=n\omega_1$ 处的频谱值；
$\omega_1\to0$ 时即连续频谱。注意周期信号的频谱密度在 $n\omega_1$ 处\textbf{聚集成冲激函数}。

写成 $F(\omega)=|F(\omega)|\mathrm{e}^{\mathrm{j}\varphi(\omega)}$，则\textbf{幅度频谱} $|F(\omega)|$
是 $\omega$ 的偶函数、\textbf{相位频谱} $\varphi(\omega)$ 是 $\omega$ 的奇函数（$f(t)$ 为实函数时）。
此处"幅频""相频"都是对频谱密度而言；实信号的频谱密度实部为 $\omega$ 的偶函数、虚部为奇函数。

\subsubsection{傅里叶反变换的三角形式}

把 $F(\omega)=|F(\omega)|\mathrm{e}^{\mathrm{j}\varphi(\omega)}$ 代入并用欧拉公式，实函数时虚部积分为零：
\begin{equation}
f(t)=\frac{1}{\pi}\int_{0}^{\infty}\big|F(\omega)\big|\cos\big(\omega t+\varphi(\omega)\big)\mathrm{d}\omega .
\end{equation}
含义：非周期信号包含了频率从零到无限大的\textbf{一切}频率的余弦分量，但各分量振幅
$\frac{1}{\pi}|F(\omega)|\mathrm{d}\omega$ 是\textbf{无穷小量}，故频谱不能再用"幅度"而要用\textbf{密度函数}表示。

\subsubsection{存在条件}

$f(t)$ 存在傅里叶变换的\textbf{充分条件}：无限区间内\textbf{绝对可积}
$\displaystyle\int_{-\infty}^{\infty}|f(t)|\mathrm{d}t<\infty$（最难满足）；
任意有限区间内只有有限个不连续点且取有限值；只有有限个极大值、极小值。
注意狄里赫利条件只是充分条件、不是必要条件——若引入冲激函数或极限处理，
不绝对可积的信号也可认为具有傅里叶变换。

\subsubsection{典型非周期信号的频谱}

\begin{center}
\begin{tabular}{lll}
\toprule
信号 & 时域表达式（$a>0$） & 频谱密度 $F(\omega)$ \\
\midrule
矩形脉冲信号 & $E\big[u(t+\frac{\tau}{2})-u(t-\frac{\tau}{2})\big]$ & $E\tau\,\mathrm{Sa}\big(\frac{\omega\tau}{2}\big)$ \\
单边指数信号 & $\mathrm{e}^{-at}u(t)$ & $\dfrac{1}{a+\mathrm{j}\omega}$ \\
双边指数信号 & $\mathrm{e}^{-a|t|}$ & $\dfrac{2a}{a^2+\omega^2}$ \\
双边奇指数信号 & $\mathrm{e}^{-at}u(t)-\mathrm{e}^{at}u(-t)$ & $-\mathrm{j}\dfrac{2\omega}{a^2+\omega^2}$ \\
单位直流信号 & $1$ & $2\pi\delta(\omega)$ \\
符号函数信号 & $\operatorname{sgn}(t)=u(t)-u(-t)$ & $\dfrac{2}{\mathrm{j}\omega}$ \\
单位冲激信号 & $\delta(t)$ & $1$ \\
单位阶跃信号 & $u(t)$ & $\pi\delta(\omega)+\dfrac{1}{\mathrm{j}\omega}$ \\
\bottomrule
\end{tabular}
\end{center}

要点（$a>0$）：
\begin{itemize}
\item \textbf{矩形脉冲}（窗函数、门函数）：$F(\omega)=\frac{2E}{\omega}\sin\frac{\omega\tau}{2}
  =E\tau\mathrm{Sa}(\frac{\omega\tau}{2})$；与周期矩形脉冲的 $F(n\omega_1)$ 相比\textbf{包络线相似、带宽相同}
  （$\frac{2\pi}{\tau}$），但一个连续、一个离散；
\item \textbf{指数类}：单边指数 $|F|=\frac{1}{\sqrt{a^2+\omega^2}}$、$\varphi=-\arctan\frac{\omega}{a}$；
  双边指数 $|F|=\frac{2a}{a^2+\omega^2}$、$\varphi=0$（实偶）；双边奇指数 $|F|=\frac{2\omega}{a^2+\omega^2}$、
  $\varphi=\pm\frac{\pi}{2}$（虚奇）；
\item \textbf{不满足狄氏条件的信号}：单位直流由双边指数取 $a\to0^+$ 极限，
  $\lim\frac{2a}{a^2+\omega^2}=0\ (\omega\neq0)$、$\infty\ (\omega=0)$，而
  $\lim 2\arctan\frac{\omega}{a}\big|_{-\infty}^{\infty}=2\pi$，故 $F(\omega)=2\pi\delta(\omega)$；
  符号函数同理取极限得 $F(\omega)=\frac{2}{\mathrm{j}\omega}$、$|F|=\frac{2}{|\omega|}$；
  单位阶跃由 $u(t)=\frac12+\frac12\operatorname{sgn}(t)$ 得 $F(\omega)=\pi\delta(\omega)+\frac{1}{\mathrm{j}\omega}$
  （含直流分量，故频谱中有冲激函数，与拉普拉斯变换不同）；
\item \textbf{单位冲激}：$F(\omega)=1$（抽样特性），频谱占据整个频率域且\textbf{均匀分布}，
  称\textbf{均匀频谱}或\textbf{白色频谱}，也可由单矩形脉冲取 $\tau\to0$ 求得；它体现\textbf{对偶性}：
  $\mathcal{F}\delta(t)=1$，$\mathcal{F}1=2\pi\delta(\omega)$，$\mathcal{F}^{-1}\delta(\omega)=\frac{1}{2\pi}$。
\end{itemize}

\subsection{傅里叶变换的性质}

\begin{center}
\begin{tabular}{lll}
\toprule
性质 & 时域 & 频域 \\
\midrule
线性 & $a_1f_1+a_2f_2$ & $a_1F_1+a_2F_2$ \\
奇偶虚实性 & $f^{*}(t)$ & $F^{*}(-\omega)$ \\
对偶性 & $F(t)$ & $2\pi f(-\omega)$ \\
尺度变换 & $f(at)$ & $\dfrac{1}{|a|}F\big(\dfrac{\omega}{a}\big)$ \\
时移 & $f(t\pm t_0)$ & $\mathrm{e}^{\pm\mathrm{j}\omega t_0}F(\omega)$ \\
频移 & $f(t)\mathrm{e}^{\pm\mathrm{j}\omega_0t}$ & $F(\omega\mp\omega_0)$ \\
微分 & $\dfrac{\mathrm{d}^nf}{\mathrm{d}t^n}$ & $(\mathrm{j}\omega)^nF(\omega)$ \\
积分 & $\displaystyle\int_{-\infty}^{t}f\mathrm{d}\tau$ & $\dfrac{F(\omega)}{\mathrm{j}\omega}+\pi F(0)\delta(\omega)$ \\
卷积（时域） & $f_1*f_2$ & $F_1F_2$ \\
卷积（频域） & $f_1f_2$ & $\dfrac{1}{2\pi}F_1*F_2$ \\
\bottomrule
\end{tabular}
\end{center}

各条要点：
\begin{itemize}
\item \textbf{线性}：傅里叶变换是线性运算，满足叠加定理；
\item \textbf{奇偶虚实性}：$f$ 为实函数时 $F(\omega)=F^{*}(-\omega)$，$|F|$ 偶、$\varphi$ 奇；
  实偶函数 $\to$ $F$ 为实偶函数（如双边指数），实奇函数 $\to$ $F$ 为虚奇函数（如双边奇指数）；
\item \textbf{对偶性}：由反变换以 $-t$ 代 $t$ 再互换 $t$ 与 $\omega$ 可证。典型体现：时域单位直流（幅度 1）
  对偶于频域冲激（强度 $2\pi$）；时域单位冲激（强度 1）对偶于频域单位直流（幅度 1）。
  例：取 $E=\frac12$、$\tau=2$ 的矩形窗有 $\mathcal{F}g(t)=\mathrm{Sa}(\omega)$，故
  $\mathcal{F}[\mathrm{Sa}(t)]=2\pi g(\omega)=\pi\ (|\omega|<1)$、$0\ (|\omega|>1)$；
\item \textbf{尺度变换}：时域压缩 $\Rightarrow$ 频带展宽（带宽增大、幅度变小），时域扩展 $\Rightarrow$
  频带压缩（带宽减小、幅度变大），频谱能量守恒；$a=-1$ 即翻转 $f(-t)\leftrightarrow F(-\omega)$。
  工程含义：压缩脉宽以提高通信速率须展宽信道频带，\textbf{通信速度与占用频带是一对矛盾}；
\item \textbf{时移}：时域延时（超前）$t_0$ 使\textbf{幅度谱不变}、相位谱产生 $-\omega t_0$（或 $+\omega t_0$）变化；
  与尺度变换组合：$\mathcal{F}[f(at-t_0)]=\frac{1}{|a|}F(\frac{\omega}{a})\mathrm{e}^{-\mathrm{j}\omega t_0/a}$；
\item \textbf{频移}：乘 $\mathrm{e}^{\pm\mathrm{j}\omega_0t}$ 使频谱右移（左移）$\omega_0$，实现\textbf{频谱搬移}；
  频移与时移在指数上标正负号的方向\textbf{相反}。\textbf{调幅信号}：
  \begin{equation}
  \mathcal{F}\big[f(t)\cos\omega_0t\big]=\tfrac12F(\omega-\omega_0)+\tfrac12F(\omega+\omega_0),\qquad
  \mathcal{F}\big[f(t)\sin\omega_0t\big]=\tfrac{\mathrm{j}}{2}F(\omega+\omega_0)-\tfrac{\mathrm{j}}{2}F(\omega-\omega_0),
  \end{equation}
  调幅、同步解调、变频都在频谱搬移的基础上完成；
\item \textbf{微分与积分}：时域微分对应频域乘 $\mathrm{j}\omega$，\textbf{增强高频}、突出信号变化
  （如 $\mathcal{F}u(t)=\pi\delta(\omega)+\frac{1}{\mathrm{j}\omega}$ 求导得 $\mathcal{F}\delta(t)=1$，
  其中 $\mathrm{j}\omega\delta(\omega)=0$ 只在 $\omega=0$ 时有意义）；时域积分对应频域乘 $\frac{1}{\mathrm{j}\omega}$，
  \textbf{弱化高频}、平滑变化，若 $F(0)=0$ 则只剩 $\frac{F(\omega)}{\mathrm{j}\omega}$
  （证明思路：把积分写成与 $u(t-\tau)$ 的卷积，用其频谱并交换积分次序）；
\item \textbf{帕斯瓦尔定理}：本课程从略，其内容与时域、频域的能量守恒一致（功率形式见 1.2.1）；
\item \textbf{卷积定理}（信号处理中应用最广的性质）：
  \begin{equation}
  \mathcal{F}\big[f_1(t)*f_2(t)\big]=F_1(\omega)F_2(\omega),\qquad
  \mathcal{F}\big[f_1(t)f_2(t)\big]=\frac{1}{2\pi}F_1(\omega)*F_2(\omega),
  \end{equation}
  证明：时域卷积定理交换积分次序并对内层用\textbf{时移特性}；频域卷积定理用反变换代入、
  交换积分次序并用频移特性。它在 LTI 系统分析中给出
  $Y(s)=U(s)G(s)$、$y(t)=u(t)*g(t)$、$Y(\mathrm{j}\omega)=U(\mathrm{j}\omega)G(\mathrm{j}\omega)$
  （$G(s)$ 传递函数，$G(\mathrm{j}\omega)$ 频率响应，$g(t)$ 单位冲激响应）。
  应用（调制与解调）：信息本身有带宽（如声音约 $22\,\mathrm{kHz}$）却无法无线传输，
  载波纯净（单频、带宽为 0）无法携带信息但高频便于传输，二者相乘即把信息频谱搬到载频附近（AM）；
  解调时通过低通滤波器检出包络、恢复原信号。
\end{itemize}

\subsection{周期信号的傅里叶变换}

动机：周期信号可用傅里叶级数（离散谱），非周期信号可用傅里叶变换（连续谱）；
通过\textbf{引入冲激函数}可得出周期信号的傅里叶变换，从而把两者的频域分析统一起来。

\subsubsection{复指数与正余弦信号的傅里叶变换}

由 $\mathcal{F}[f(t)\mathrm{e}^{\mathrm{j}\omega_0t}]=F(\omega-\omega_0)$ 并令 $f(t)=1$、
$F(\omega)=2\pi\delta(\omega)$，得
\begin{equation}
\mathcal{F}\big[\mathrm{e}^{\mathrm{j}\omega_0t}\big]=2\pi\delta(\omega-\omega_0),
\end{equation}
\begin{equation}
\mathcal{F}\big[\sin\omega_0t\big]=-\mathrm{j}\pi\delta(\omega-\omega_0)+\mathrm{j}\pi\delta(\omega+\omega_0),\qquad
\mathcal{F}\big[\cos\omega_0t\big]=\pi\delta(\omega-\omega_0)+\pi\delta(\omega+\omega_0).
\end{equation}

\subsubsection{一般周期信号}

设 $f(t)=\sum F_n\mathrm{e}^{\mathrm{j}n\omega_1t}$，由 $\mathcal{F}[\mathrm{e}^{\mathrm{j}n\omega_1t}]
=2\pi\delta(\omega-n\omega_1)$ 得
\begin{equation}
F(\omega)=\sum_{n=-\infty}^{\infty}2\pi F_n\delta(\omega-n\omega_1),
\qquad
F_n=\frac{1}{T_1}\int_{-\frac{T_1}{2}}^{\frac{T_1}{2}}f(t)\mathrm{e}^{-\mathrm{j}n\omega_1t}\mathrm{d}t .
\end{equation}
冲激位于各\textbf{谐频}（$0,\pm\omega_1,\pm2\omega_1,\dots$）处，强度为 $F_n$ 的 $2\pi$ 倍。
结论：\textbf{周期信号的傅里叶变换（频谱密度函数）由无穷多个冲激函数组成}。

\subsubsection{周期信号傅里叶系数与单脉冲傅里叶变换的关系}

设单脉冲信号的傅里叶变换为
$F_0(\omega)=\displaystyle\int_{-T_1/2}^{T_1/2}f(t)\mathrm{e}^{-\mathrm{j}\omega t}\mathrm{d}t$，则
\begin{equation}
F_n=\frac{1}{T_1}F_0(\omega)\Big|_{\omega=n\omega_1},
\end{equation}
据此可由单脉冲的傅里叶变换很方便地求出周期脉冲序列的傅里叶系数。

\subsubsection{周期矩形脉冲信号与周期单位冲激序列}

由 $F_n=\frac{E\tau}{T_1}\mathrm{Sa}(\frac{n\omega_1\tau}{2})$ 得
\begin{equation}
F(\omega)=\sum_{n=-\infty}^{\infty}2\pi\frac{E\tau}{T_1}\mathrm{Sa}\Big(\frac{n\omega_1\tau}{2}\Big)\delta(\omega-n\omega_1)
=E\tau\omega_1\sum_{n=-\infty}^{\infty}\mathrm{Sa}\Big(\frac{n\omega_1\tau}{2}\Big)\delta(\omega-n\omega_1),
\end{equation}
它只在 $\omega=n\omega_1$ 时有意义；与傅里叶级数系数相比，二者有\textbf{相同的包络线}、都是离散的，
一个表现为离散函数、一个表现为\textbf{冲激序列}。

设 $\delta_T(t)=\sum_{n=-\infty}^{\infty}\delta(t-nT_1)$，由 $\mathcal{F}\delta(t)=1$ 及
$F_n=\frac{1}{T_1}\mathcal{F}[\delta(t)]|_{\omega=n\omega_1}=\frac{1}{T_1}$ 得
\begin{equation}
\delta_T(t)=\frac{1}{T_1}\sum_{n=-\infty}^{\infty}\mathrm{e}^{\mathrm{j}n\omega_1t},\qquad
\mathcal{F}\big[\delta_T(t)\big]=\sum_{n=-\infty}^{\infty}\frac{2\pi}{T_1}\delta(\omega-n\omega_1)
=\sum_{n=-\infty}^{\infty}\omega_1\delta(\omega-n\omega_1),
\end{equation}
即周期单位冲激序列的傅里叶变换是间隔 $\omega_1$、强度 $\omega_1$ 的等强度冲激序列。

\section{连续信号的复频域分析}

\subsection{从傅里叶变换到拉普拉斯变换}

傅里叶变换要求信号满足狄里赫利条件，而：直流信号、周期信号等不满足该条件，需引入冲激函数或极限求解；
功率型非周期信号、指数增长型信号 $\mathrm{e}^{at}$（$a>0$）等难以求解傅里叶变换。
\textbf{拉普拉斯变换}（拉氏变换，Laplace Transform）把频域\textbf{推广}到复频域，
克服傅里叶变换的局限性，扩大频谱分析的范围。

做法：对不满足绝对可积的信号（如 $\mathrm{e}^{at}$，$a>0$），先乘随时间衰减的因子
$\mathrm{e}^{-\sigma t}$（$\sigma>0$，使 $\lim_{t\to\infty}\mathrm{e}^{at}\mathrm{e}^{-\sigma t}=0$），
使其满足绝对可积条件；令 $s=\sigma+\mathrm{j}\omega$（$\sigma\in\mathbb{R}$），得
\begin{equation}
F_b(s)=\int_{-\infty}^{\infty}f(t)\mathrm{e}^{-st}\,\mathrm{d}t .
\end{equation}
反变换由傅里叶反变换式两边同乘 $\mathrm{e}^{\sigma t}$、用 $\mathrm{d}s=\mathrm{j}\mathrm{d}\omega$、
$\omega\to\pm\infty$ 对应 $s\to\sigma\pm\mathrm{j}\infty$ 得到
\begin{equation}
f(t)=\frac{1}{2\pi\mathrm{j}}\int_{\sigma-\mathrm{j}\infty}^{\sigma+\mathrm{j}\infty}F_b(s)\mathrm{e}^{st}\,\mathrm{d}s .
\end{equation}

\subsection{拉普拉斯变换对}

\textbf{双边拉普拉斯变换对}（积分上下限包括时域正、负区间）：
\begin{equation}
F_b(s)=\int_{-\infty}^{\infty}f(t)\mathrm{e}^{-st}\mathrm{d}t,\qquad
f(t)=\frac{1}{2\pi\mathrm{j}}\int_{\sigma-\mathrm{j}\infty}^{\sigma+\mathrm{j}\infty}F_b(s)\mathrm{e}^{st}\mathrm{d}s,
\qquad f(t)\overset{\mathcal{L}}{\longleftrightarrow}F_b(s).
\end{equation}
它表明信号 $f(t)$ 是复指数信号 $\mathrm{e}^{st}$ 的线性组合：$\sigma\neq0$ 时为增幅振荡（$\sigma>0$）
或减幅振荡（$\sigma<0$）；$\sigma=0$ 时退化为傅里叶变换
$F_b(\mathrm{j}\omega)=\int f(t)\mathrm{e}^{-\mathrm{j}\omega t}\mathrm{d}t=F(\omega)$。

\textbf{单边拉普拉斯变换}：当 $f(t)=0,\ \forall t<0$ 时双边与单边一致，且
$\mathcal{L}[f(t)]=\mathcal{F}[f(t)u(t)\mathrm{e}^{-\sigma t}]\big|_{s=\sigma+\mathrm{j}\omega}$。
关系：双边拉氏变换在 $\sigma=0$ 处退化为傅里叶变换；有始信号的双边拉氏变换即单边拉氏变换。

\subsection{收敛域与收敛坐标}

为使 $\mathrm{e}^{at}$ 乘衰减因子后收敛，须 $\sigma>a$，否则 $t\to\infty$ 时仍不收敛。一般要求
\begin{equation}
\lim_{t\to\infty}f(t)\mathrm{e}^{-\sigma t}=0,\qquad \sigma>\sigma_0,
\end{equation}
此时 $f(t)\mathrm{e}^{-\sigma t}$ 在 $\sigma>\sigma_0$ 上收敛、积分存在，可作拉普拉斯变换。
满足上式的函数称"\textbf{指数阶函数}"；若发散，可借助指数函数的衰减变为收敛函数。
$\sigma_0$ 与 $f(t)$ 的性质有关，称为\textbf{收敛坐标}，按 $\sigma_0$ 在 $s$ 平面上划分的区域即\textbf{收敛域}。
拉普拉斯变换是傅里叶变换的\textbf{推广}：存在傅里叶变换的信号一般也存在拉普拉斯变换，
但个别除外（如 $f(t)=1$，$F(\omega)=2\pi\delta(\omega)$，却不存在拉氏变换）；
存在拉普拉斯变换的信号，其傅里叶变换不一定存在。

\subsection{由拉普拉斯变换求傅里叶变换}

首要条件：$f(t)$ 为\textbf{有始信号}（$f(t)=0,\ \forall t<0$）。按收敛边界 $\sigma_0$ 分三种情形：

\begin{center}
\begin{tabular}{ll}
\toprule
收敛边界 & 傅里叶变换 \\
\midrule
$\sigma_0>0$ & 不存在 \\
$\sigma_0<0$ & 令 $s=\mathrm{j}\omega$，$F(\omega)=F_b(s)|_{s=\mathrm{j}\omega}$ \\
$\sigma_0=0$（边界位于虚轴） & $F_b(s)|_{s=\mathrm{j}\omega}+\sum K_n\pi\delta(\omega-\omega_n)$ \\
\bottomrule
\end{tabular}
\end{center}

当 $\sigma_0=0$ 且虚轴上有 $N$ 个极点时，
\begin{equation}
F_b(s)=F_{b1}(s)+\sum_{n=1}^{N}\frac{K_n}{s-\mathrm{j}\omega_n},
\qquad
\mathcal{F}\big[f(t)\big]=F_b(s)\Big|_{s=\mathrm{j}\omega}+\sum_{n=1}^{N}K_n\pi\delta(\omega-\omega_n),
\end{equation}
其中 $F_{b1}(s)$ 由位于 $s$ 左半平面极点对应的分式构成：第一部分直接通过 $s=\mathrm{j}\omega$ 求解，
第二部分由虚轴上每个极点 $\mathrm{j}\omega_n$ 对应的冲激项组成。

\subsection{本章变换类型小结}

\begin{center}
\begin{tabular}{lll}
\toprule
信号类别 & 变换类型 & 适用对象 \\
\midrule
连续信号（本章） & 傅里叶变换（FT） & 连续非周期信号 \\
 & 傅里叶级数（FS） & 连续周期信号 \\
离散信号（第二章） & 离散时间傅里叶变换（DTFT） & 离散非周期信号 \\
 & 离散傅里叶变换（DFT） & 离散周期信号 \\
\bottomrule
\end{tabular}
\end{center}

\chapter{离散信号的分析}
\providecommand{\Saa}{\operatorname{Sa}}

\section{信号的采样和恢复}

\subsection{离散信号与采样信号}

\begin{itemize}
\item \textbf{离散信号}：只在某些不连续的规定时刻给出瞬时值；可由连续信号采样得到，也可能本身就是离散的（如逐年人口）。
\item \textbf{采样信号}：用采样脉冲 $p(t)$ 从连续信号中“抽取”一系列离散样值，记 $f_s(t)$（$s$ 对应 sampling）。
\item 采样信号\textbf{时间离散、幅值连续}；再经量化、编码才成为数字信号。
\end{itemize}
采样要回答两点：$f_s(t)$ 的频谱与 $f(t)$ 频谱的关系；$f_s(t)$ 能否无失真地恢复 $f(t)$。

\subsection{采样信号的频谱}

由频域卷积定理（时域相乘对应频域卷积），$f_s(t)=f(t)p(t)$ 的频谱为
\begin{equation}
F_s(\omega)=\frac{1}{2\pi}F(\omega)*P(\omega).
\end{equation}
周期 $T_s$ 的采样脉冲可展成 $p(t)=\sum_n P_n e^{jn\omega_s t}$，
故 $P(\omega)=2\pi\sum_n P_n\delta(\omega-n\omega_s)$，其中
\begin{equation}
P_n=\frac{1}{T_s}\int_{-T_s/2}^{T_s/2}p(t)e^{-jn\omega_s t}\,\mathrm{d}t,
\qquad \omega_s=\frac{2\pi}{T_s}.
\end{equation}
利用 $F(\omega)*\delta(\omega-n\omega_s)=F(\omega-n\omega_s)$ 得\textbf{一般式}
\begin{equation}
F_s(\omega)=\sum_{n=-\infty}^{\infty}P_n F(\omega-n\omega_s),
\end{equation}
即采样使原频谱\textbf{以 $\omega_s$ 为周期搬移并叠加}，权重为各次谐波系数 $P_n$。

\subsubsection{矩形脉冲采样}
$p(t)$ 为周期 $T_s$、幅度 $E$、脉宽 $\tau\ll T_s$ 的矩形脉冲时
\begin{equation}
P_n=\frac{1}{T_s}\int_{-\tau/2}^{\tau/2}E e^{-jn\omega_s t}\,\mathrm{d}t
=\frac{E\tau}{T_s}\,\Saa\!\left(\frac{n\omega_s\tau}{2}\right),
\end{equation}
\begin{equation}
F_s(\omega)=\frac{E\tau}{T_s}\sum_{n=-\infty}^{\infty}
\Saa\!\left(\frac{n\omega_s\tau}{2}\right)F(\omega-n\omega_s),
\end{equation}
各次搬移的幅度按包络 $\Saa(n\omega_s\tau/2)$ 衰减，包络在 $\omega=2k\pi/\tau$ 处过零。
$\tau\to0$、$E\tau\to1$ 时退化为冲激序列，包络趋于常数。

\subsubsection{冲激采样（理想采样）}
取 $p(t)=\delta_{T_s}(t)=\sum_n\delta(t-nT_s)$，则
\begin{equation}
f_s(t)=f(t)\sum_{n=-\infty}^{\infty}\delta(t-nT_s)
=\sum_{n=-\infty}^{\infty}f(nT_s)\,\delta(t-nT_s),
\end{equation}
由 $P_n=1/T_s$ 得
\begin{equation}
F_s(\omega)=\frac{1}{T_s}\sum_{n=-\infty}^{\infty}F(\omega-n\omega_s).
\end{equation}
即冲激采样后原频谱\textbf{等幅}地以 $\omega_s$ 为周期重复，幅度乘以常数 $1/T_s$。

\begin{center}
\begin{tabular}{lll}
\toprule
项目 & 矩形脉冲采样 & 冲激（理想）采样 \\
\midrule
$p(t)$ & 周期矩形，幅度 $E$、脉宽 $\tau$ & $\delta_{T_s}(t)=\sum_n\delta(t-nT_s)$ \\
$P_n$ & $\dfrac{E\tau}{T_s}\Saa\!\left(\dfrac{n\omega_s\tau}{2}\right)$ & $\dfrac{1}{T_s}$ \\
$F_s(\omega)$ & 按 $\Saa$ 包络衰减的搬移叠加 & 等幅搬移叠加 \\
\bottomrule
\end{tabular}
\end{center}

\subsection{时域采样定理与频谱混叠}

\textbf{时域采样定理（香农定理）}：若 $f(t)$ 频谱受限，只占据 $-\omega_m\sim+\omega_m$
（最高频率 $f_m=\omega_m/2\pi$），则 $f(t)$ 可由等间隔采样值唯一确定，只要
\begin{equation}
T_s\le\frac{\pi}{\omega_m}=\frac{1}{2f_m},
\qquad\text{即}\qquad \omega_s\ge2\omega_m\ (f_s\ge2f_m).
\end{equation}
$\omega_s=2\omega_m$ 称\textbf{奈奎斯特频率}，$T_s=\pi/\omega_m=1/(2f_m)$ 称\textbf{奈奎斯特间隔}。
定理要求三点：信号带限、等间隔采样、采样率不低于 $2f_m$。

\textbf{频谱混叠}：$\omega_s<2\omega_m$ 时相邻副本 $F(\omega-n\omega_s)$ 重叠，高频被折叠到低频段。
混叠后各分量已叠加在一起，\textbf{无法再用滤波器分开}，属不可逆失真。

\textbf{易错点}
\begin{itemize}
\item $2f_m$ 是临界值，工程上须留裕量 $\omega_s>2\omega_m$。
\item $f_m$（频率）与 $\omega_m$（角频率）相差 $2\pi$，条件不能混写。
\item 采样后频谱幅度多出 $1/T_s$，比较幅度时不能漏。
\end{itemize}

\subsection{频域采样定理}

对连续频谱以间隔 $\omega_0$ 采样，得 $F_p(\omega)=F(\omega)\sum_k\delta(\omega-k\omega_0)$。
由时域卷积定理
\begin{equation}
f_p(t)=f(t)*\frac{1}{\omega_0}\sum_{k=-\infty}^{\infty}\delta(t-kT_0)
=\frac{1}{\omega_0}\sum_{k=-\infty}^{\infty}f(t-kT_0),\qquad T_0=\frac{2\pi}{\omega_0},
\end{equation}
即频域采样对应\textbf{时域以 $T_0$ 为周期的重复}，幅度乘 $1/\omega_0$。

\textbf{频域采样定理}：若 $f(t)$ 时间受限，集中在 $-t_m\sim+t_m$，则以不大于 $1/(2t_m)$ 的频率间隔
对 $F(\omega)$ 采样，可由 $F_p(\omega)$ 唯一恢复 $F(\omega)$（要求 $T_0\ge2t_m$，时域副本不重叠）。

\begin{center}
\begin{tabular}{lll}
\toprule
 & 时域采样 & 频域采样 \\
\midrule
条件 & 频谱受限（带限） & 时间受限（时限） \\
间隔要求 & $T_s\le1/(2f_m)$ & $\omega_0\le1/(2t_m)$ \\
另一域变化 & 频谱以 $\omega_s$ 周期重复 & 时域以 $T_0$ 周期重复 \\
恢复手段 & 频域加理想低通 & 时域乘矩形窗 \\
\bottomrule
\end{tabular}
\end{center}
一般规律：\textbf{时域离散化 $\Leftrightarrow$ 频域周期化；时域周期化 $\Leftrightarrow$ 频域离散化}。

\subsection{采样信号的恢复与内插公式}

在频域取主频带，令
\begin{equation}
H(\omega)=\begin{cases}T_s, & |\omega|\le\omega_c,\\ 0, & |\omega|>\omega_c,\end{cases}
\qquad \omega_c=\omega_m,\ \omega_s=2\omega_m,
\end{equation}
则其冲激响应为
\begin{equation}
h(t)=\Saa\!\left(\frac{\omega_s t}{2}\right),
\end{equation}
由时域卷积定理得\textbf{内插公式}
\begin{equation}
f(t)=f_s(t)*h(t)=\sum_{n=-\infty}^{\infty}f(nT_s)\,\Saa\!\left[\frac{\omega_s}{2}(t-nT_s)\right],
\end{equation}
即带限信号任一时刻的值由全部采样值经取样函数加权叠加确定；内插函数在各采样点为 $1$、其余采样点为 $0$（无码间串扰）。
对偶的频域内插公式为
\begin{equation}
F(\omega)=\frac{1}{2\pi}F_p(\omega)*G(\omega)
=\sum_{k=-\infty}^{\infty}F(k\omega_0)\,\Saa\!\left[\frac{T_0}{2}(\omega-k\omega_0)\right].
\end{equation}
理想低通非因果、无限长，不可实现，故实际只能近似恢复。

\textbf{采样频率的选择}：一个周期至少采样两点才能保留该频率分量信息；快变信号要提高采样率，
过高则增大存储量并使采样不稳；非带限信号采样前须加\textbf{抗混叠滤波器}滤除次要高频；
工程上常取 $f_s>10f_m$。实际采样与滤波均非理想，抗混叠滤波、A/D、D/A 各环节都会引入误差。

\section{离散信号的时域描述和运算}

\subsection{序列及其表示}

等间隔时刻取值即可用序列 $x(n)$ 表示：$n$ 为整数序号，$x(n)$ 为第 $n$ 点的\textbf{样值}。表示方法有
\textbf{列举法}（$\{x(n)\}=\{\cdots,0,1,2,3,4,3,2,1,0,\cdots\}$，标出 $n=0$ 位置）、
\textbf{闭式法}（如 $x(n)=4-n\ (n\le3)$）、\textbf{图形法}（$n$ 为横轴的谱线图，可画包络线）。
离散信号能量定义为
\begin{equation}
W=\sum_{n=-\infty}^{\infty}\big|x(n)\big|^2 .
\end{equation}

\subsection{典型离散信号（序列）}

\begin{enumerate}
\item \textbf{单位脉冲序列}
\begin{equation}
\delta(n)=\begin{cases}1, & n=0,\\ 0, & n\ne0,\end{cases}
\qquad \sum_{n=-\infty}^{\infty}\delta(n)=1 .
\end{equation}
取样特性：$x(n)\delta(n)=x(0)\delta(n)$，
$\sum\limits_{n=-\infty}^{\infty}x(n)\delta(n-n_0)=x(n_0)$。

\item \textbf{单位阶跃序列}
\begin{equation}
u(n)=\begin{cases}1, & n\ge0,\\ 0, & n<0,\end{cases}
\qquad \delta(n)=u(n)-u(n-1),\quad u(n)=\sum_{k=0}^{\infty}\delta(n-k).
\end{equation}
注意 $u(0)=1$，与 $u(t)$ 在 $t=0$ 处跳变、取值不确定不同。

\item \textbf{矩形序列}
\begin{equation}
R_N(n)=\begin{cases}1, & 0\le n\le N-1,\\ 0, & \text{其他},\end{cases}
\qquad R_N(n)=u(n)-u(n-N)=\sum_{k=0}^{N-1}\delta(n-k).
\end{equation}

\item \textbf{斜变序列}：$x(n)=n\,u(n)$。

\item \textbf{实指数序列}：$x(n)=a^n u(n)$。$|a|<1$ 收敛，$|a|>1$ 发散；$a>0$ 同号，$a<0$ 正负摆动。

\item \textbf{正弦型序列}：
\begin{equation}
x(n)=A\sin(\Omega_0 n+\phi),
\end{equation}
$\Omega_0$ 为\textbf{数字角频率}（rad），$\phi$ 为初相。若由 $A\sin(\omega_0t+\phi)$ 以 $T_s$ 采样得到，
则 $\Omega_0=\omega_0T_s=\omega_0/f_s$。

\item \textbf{复指数序列}：
\begin{equation}
x(n)=e^{(\sigma+j\Omega_0)n}=e^{\sigma n}\big(\cos\Omega_0 n+j\sin\Omega_0 n\big),
\end{equation}
可统一表示实指数、正弦、直流序列。

\item \textbf{任意序列}：
\begin{equation}
x(n)=\sum_{m=-\infty}^{\infty}x(m)\,\delta(n-m).
\end{equation}
\end{enumerate}

\begin{center}
\begin{tabular}{lll}
\toprule
对比项 & 连续 $\delta(t)$ & 离散 $\delta(n)$ \\
\midrule
非零处取值 & $t=0$ 处为 $\infty$（奇异函数） & $n=0$ 处为有限值 $1$ \\
积分/求和 & $\int_{-\infty}^{\infty}\delta(t)\mathrm{d}t=1$ & $\sum_{n=-\infty}^{\infty}\delta(n)=1$ \\
与阶跃关系 & $\delta(t)=\mathrm{d}u(t)/\mathrm{d}t$ & $\delta(n)=u(n)-u(n-1)$ \\
\bottomrule
\end{tabular}
\end{center}

\subsection{正弦型序列与复指数序列的特性}

\textbf{（1）离散正弦序列不一定是周期序列。}要求 $x(n+N)=x(n)$，即 $N\Omega_0=2k\pi$，故
\begin{equation}
N=\frac{2k\pi}{\Omega_0},\qquad k=\pm1,\pm2,\cdots
\end{equation}
只有 $\dfrac{2\pi}{\Omega_0}$ 为\textbf{有理数}时，取使 $N$ 为最小正整数的 $k$ 才周期；
若为无理数则\textbf{非周期}（包络仍是周期正弦）。
连续正弦恒周期（$T=2\pi/\omega_0$）；采样后若一个连续周期内样点数非整数，序列周期会变长甚至不存在。

\textbf{（2）时移与相移不一定对应。}连续 $A\sin(\omega_0t+\varphi)$ 中时移与相移一一对应；
离散 $A\sin(\Omega_0n+\phi)$ 中时移 $n\pm n_0$ 必对应相移 $\mp\Omega_0n_0$，但给定相移 $\phi_0$，
只有它能写成 $\Omega_0n_0$（$n_0$ 为整数，即 $\phi_0$ 是 $\Omega_0$ 的整数倍）时才对应整数位移。

\textbf{（3）数字频率的模糊性。}因 $e^{j(\Omega\pm2k\pi)n}=e^{j\Omega n}$，故
\begin{equation}
\Omega_1=\Omega_2+2k\pi\ \Longrightarrow\ x_1(n)=x_2(n),
\end{equation}
不同的数字频率可对应完全相同的序列。故 $\Omega$ 只需在长 $2\pi$ 的区间内取值，通常取
\begin{equation}
\Omega\in[0,2\pi)\quad\text{或}\quad \Omega\in[-\pi,\pi).
\end{equation}

\subsection{序列的时域运算}

\begin{itemize}
\item \textbf{平移} $x(n-m)$：$m>0$ 右移（延时），$m<0$ 左移；$n,m$ 均整数。\textbf{翻转} $x(-n)$：以 $n=0$ 纵轴为对称轴。
\item \textbf{和、积}：$z(n)=x(n)\pm y(n)$、$z(n)=x(n)y(n)$，同序号逐项运算。
\item \textbf{累加}：
\begin{equation}
y(n)=\sum_{k=-\infty}^{n}x(k),
\end{equation}
即 $n$ 及 $n$ 以前全部样值之和（对应连续积分）。
\item \textbf{差分}：
\begin{equation}
\Delta x(n)=x(n+1)-x(n),\qquad
\nabla x(n)=x(n)-x(n-1)=\Delta x(n-1).
\end{equation}
\item \textbf{时间尺度变换} $x(mn)$、$x(n/m)$：$x(2n)$ 是每隔 2 点取 1 点（抽取，采样率降为 $1/m$），
并非波形压缩一半；$x(n/2)$ 在相邻样点间插入新样点（插值，间隔变为 $T_s/2$），插入后序号需重排。
\end{itemize}

\subsection{离散系统的单位脉冲响应与卷积和}

设离散线性系统对 $\delta(n-m)$ 的响应为 $h_m(n)$，由线性叠加，输入 $x(n)=\sum_m x(m)\delta(n-m)$ 的响应为
\begin{equation}
y(n)=\sum_{m=-\infty}^{\infty}x(m)h_m(n).
\end{equation}
若系统又时不变，则 $h_m(n)=h_0(n-m)$，记 $h(n)=h_0(n)$ 为\textbf{单位脉冲响应（单位样值响应）}，于是
\begin{equation}
y(n)=\sum_{m=-\infty}^{\infty}x(m)h(n-m)=x(n)*h(n),
\end{equation}
称为\textbf{卷积和（叠加和）}。它表明 \textbf{LTI 系统的 $h(n)$ 完全刻画了系统特性}。

\subsection{卷积和的计算与性质}

\textbf{计算步骤（图解法）}
\begin{enumerate}
\item \textbf{换坐标}：$n\to m$，把 $n$ 视为参变量；
\item \textbf{翻转}：将 $h(m)$ 关于 $m=0$ 翻转为 $h(-m)$；
\item \textbf{平移}：对某定值 $n$，把 $h(-m)$ 平移 $n$ 得 $h(n-m)$；
\item \textbf{相乘}：$h(n-m)$ 与 $x(m)$ 在相同 $m$ 处逐点相乘；
\item \textbf{累加}：乘积相加得该 $n$ 处的 $y(n)$，对各 $n$ 重复。
\end{enumerate}
$h(n-m)$ 是“先翻转、后移位”，$n$ 增大时向右移；非零区间不重合时 $y(n)=0$。

\textbf{主要性质}
\begin{align}
&\text{交换律：} && x(n)*h(n)=h(n)*x(n),\\
&\text{分配律：} && x(n)*\big[h_1(n)+h_2(n)\big]=x(n)*h_1(n)+x(n)*h_2(n),\\
&\text{结合律：} && \big[x(n)*h_1(n)\big]*h_2(n)=x(n)*\big[h_1(n)*h_2(n)\big],\\
&\text{与脉冲卷积：} && x(n-m)*\delta(n-k)=x(n-m-k).
\end{align}
最后一条是化简关键：把 $x(n)$、$h(n)$ 写成 $\sum a_k\delta(n-k)$，卷积和即两组系数的“多项式乘法”。
若 $x(n)$ 长 $N_1$、$h(n)$ 长 $N_2$，则 $y(n)$ 长度为
\begin{equation}
N_1+N_2-1 .
\end{equation}

\section{离散信号的频域分析}

离散信号频域分析要解决两件事：信号离散化后频谱发生了什么变化（由 DFS 推出）；
如何用计算机快速计算（DFT 与 FFT）。各变换对应关系：
连续周期—CFS，连续非周期—CTFT 与拉普拉斯变换；离散周期—DFS，
离散非周期—DTFT、DFT、FFT 与 Z 变换。

\subsection{离散傅里叶级数（DFS）}

连续复指数 $e^{j\omega t}$ 的频率 $\omega\in(-\infty,\infty)$；而离散复指数满足
\begin{equation}
e^{j(\Omega\pm2k\pi)n}=e^{j\Omega n},
\end{equation}
故它作为 $\Omega$ 的函数以 $2\pi$ 为周期，无限频率范围被映射到长度 $2\pi$ 的区间内——
这是“时域离散化 $\Rightarrow$ 频域周期化”的根源。

\textbf{推导}：对周期 $T_0$ 的连续信号作 $N$ 点采样，$T_0=NT_s$，则
\begin{equation}
\omega_0=\frac{2\pi}{T_0}=\frac{2\pi}{NT_s},\qquad
\Omega_0=\omega_0T_s=\frac{2\pi}{N},
\end{equation}
$\Omega_0$ 为基本数字频率，$k\Omega_0$ 为 $k$ 次谐波数字频率。把连续傅里叶系数中的积分换成一个周期内的累加
（$\mathrm{d}t\to T_s$），得 $X(k\Omega_0)=\frac{1}{N}\sum_{n=0}^{N-1}x(n)e^{-jk\Omega_0n}$。
由于 $e^{jk\Omega_0n}$ 只在 $k$ 的 $N$ 个相继取值上互异，$x(n)$ 只需 $N$ 项即可展开：
\begin{align}
x(n)&=\sum_{k=\langle N\rangle}X(k\Omega_0)\,e^{jk\Omega_0 n},\\
X(k\Omega_0)&=\frac{1}{N}\sum_{n=\langle N\rangle}x(n)\,e^{-jk\Omega_0 n},
\end{align}
$\sum\limits_{n=\langle N\rangle}$ 表示在任一个周期（$N$ 点）内求和。

\textbf{要点}
\begin{itemize}
\item $X(k\Omega_0)$ 为\textbf{频谱系数}（复振幅），周期序列被分解为 $N$ 个谐波关系的复指数之和。
\item $X\big((k+N)\Omega_0\big)=X(k\Omega_0)$，频谱是以 $N$ 为周期的离散序列，取任一周期即可。
\item 离散周期序列 $\leftrightarrow$ 周期离散频谱，一一对应。
\end{itemize}

\textbf{与连续周期信号频谱的关系}：若 $x(n)$ 由连续周期信号采样得到（每周期 $N$ 点），
则 $X(k\Omega_0)$ 可看作 $X(k\omega_0)$ 的周期延拓（周期 $\omega_s$ 或数字频率 $2\pi$）并在频率轴乘 $T_s$；
一个周期内两者幅频数值相等、频率标度不同。满足采样定理时可由一个周期恢复 $X(k\omega_0)$，否则出现混叠、只能近似。

\subsection{DFS 的主要性质}

设 $x(n)\xleftrightarrow{\ \mathrm{DFS}\ }X(k\Omega_0)$，$h(n)\xleftrightarrow{\ \mathrm{DFS}\ }H(k\Omega_0)$，周期均 $N$。

\begin{enumerate}
\item \textbf{线性}：$ax(n)+by(n)\xleftrightarrow{\ \mathrm{DFS}\ }aX(k\Omega_0)+bY(k\Omega_0)$。
\item \textbf{周期卷积定理}：
\begin{equation}
x(n)\circledast h(n)\ \xleftrightarrow{\ \mathrm{DFS}\ }\ N\,X(k\Omega_0)H(k\Omega_0),
\qquad
x(n)h(n)\ \xleftrightarrow{\ \mathrm{DFS}\ }\ X(k\Omega_0)\circledast H(k\Omega_0),
\end{equation}
其中周期卷积只在\textbf{一个周期内}求和：
\begin{equation}
x(n)\circledast h(n)=\sum_{m=0}^{N-1}x(m)h(n-m)=h(n)\circledast x(n).
\end{equation}
\item \textbf{复共轭}：$x^*(-n)\xleftrightarrow{\ \mathrm{DFS}\ }X^*(k\Omega_0)$。
\item \textbf{位移}：
\begin{equation}
x(n-m)\ \xleftrightarrow{\ \mathrm{DFS}\ }\ e^{-jk\Omega_0 m}X(k\Omega_0).
\end{equation}
\item \textbf{帕斯瓦尔定理}：
\begin{equation}
\frac{1}{N}\sum_{n=0}^{N-1}\big|x(n)\big|^2=\sum_{k=0}^{N-1}\big|X(k\Omega_0)\big|^2 .
\end{equation}
\end{enumerate}
\textbf{易错点}：DFS 的 $1/N$ 在\textbf{正变换}，故周期卷积定理带因子 $N$；DFT 的 $1/N$ 在反变换，
同一性质写成 $X(k)H(k)$ 不带 $N$。各教材约定不同，\textbf{必须与所用定义配套}。

\subsection{离散时间傅里叶变换（DTFT）}

把长 $N_1+N_2$ 的有限长序列以 $N>N_1+N_2$ 为周期延拓成 $x_N(n)=\sum_m x(n-mN)$，再令 $N\to\infty$：
周期序列回到非周期序列，DFS 求和变为无穷级数，而 $X(k\Omega_0)=\frac{1}{N}X(e^{jk\Omega_0})$ 趋于无穷小。
于是定义\textbf{频谱密度函数}
\begin{align}
X(\Omega)=X(e^{j\Omega})&=\sum_{n=-\infty}^{\infty}x(n)e^{-j\Omega n},\\
x(n)&=\frac{1}{2\pi}\int_{-\pi}^{\pi}X(\Omega)e^{j\Omega n}\,\mathrm{d}\Omega ,
\end{align}
称\textbf{离散时间傅里叶变换（DTFT）}。$X(\Omega)$ 以 $2\pi$ 为周期（$e^{-j2\pi n}=1$），
是 $\Omega$ 的\textbf{连续周期}函数，表示频谱密度而非离散谱线。

\textbf{存在条件}：无限长序列要求绝对可和
\begin{equation}
\sum_{n=-\infty}^{\infty}\big|x(n)\big|<\infty,
\end{equation}
或能量有限 $\sum_n|x(n)|^2<\infty$，则级数一致收敛。

\textbf{两个基本变换对}（几何级数求和）
\begin{align}
x(n)&=a^n u(n),\ |a|<1: & X(\Omega)&=\frac{1}{1-ae^{-j\Omega}},\\
x(n)&=-a^n u(-n-1),\ |a|>1: & X(\Omega)&=\frac{1}{1-ae^{-j\Omega}}.
\end{align}
两式形式相同但收敛域不同（$|a|<1$ 与 $|a|>1$），说明
\textbf{只看 $X(\Omega)$ 的表达式不能唯一确定序列}。

\textbf{有限长矩形序列}：$x(n)=1$（$-N_1\le n\le N_1$）时
\begin{equation}
X(\Omega)=\frac{\sin\!\big[(N_1+\tfrac12)\Omega\big]}{\sin(\Omega/2)},
\end{equation}
频谱连续、以 $2\pi$ 为周期，起伏由截断引起（对应连续矩形脉冲的 $\Saa$ 谱，但这里周期化）。

\subsection{DTFT 的性质}

设 $x(n)\xleftrightarrow{\ \mathrm{DTFT}\ }X(\Omega)$，$y(n)\xleftrightarrow{\ \mathrm{DTFT}\ }Y(\Omega)$。

\begin{center}
\begin{tabular}{lll}
\toprule
性质 & 时域 & 频域 \\
\midrule
线性 & $ax(n)+by(n)$ & $aX(\Omega)+bY(\Omega)$ \\
时域平移 & $x(n-n_0)$ & $e^{-j\Omega n_0}X(\Omega)$ \\
时间翻转 & $x(-n)$ & $X(-\Omega)$ \\
频域平移 & $e^{j\Omega_0 n}x(n)$ & $X(\Omega-\Omega_0)$ \\
时域卷积 & $x(n)*y(n)$ & $X(\Omega)Y(\Omega)$ \\
共轭对称 & $x^*(n)$ & $X^*(-\Omega)$ \\
频域微分 & $n\,x(n)$ & $j\dfrac{\mathrm{d}X(\Omega)}{\mathrm{d}\Omega}$ \\
频域卷积（相乘） & $x(n)y(n)$ & $\dfrac{1}{2\pi}\displaystyle\int_{-\pi}^{\pi}X(\lambda)Y(\Omega-\lambda)\mathrm{d}\lambda$ \\
\bottomrule
\end{tabular}
\end{center}

\textbf{用 DTFT 解差分方程（零初始条件）}
\begin{enumerate}
\item 两边取 DTFT，用时移性质把 $y(n-k)$ 化为 $e^{-j\Omega k}Y(\Omega)$；
\item 解出 $Y(\Omega)=H(\Omega)X(\Omega)$，
\begin{equation}
H(\Omega)=\frac{Y(\Omega)}{X(\Omega)}=\frac{\sum\limits_{k=0}^{M}b_k e^{-j\Omega k}}{1+\sum\limits_{k=1}^{N}a_k e^{-j\Omega k}};
\end{equation}
\item 把 $H(\Omega)$ 拆成基本变换对可表示的部分分式，逐项反变换得 $y(n)$。
\end{enumerate}

\textbf{与 DFS、CTFT 的关系}
\begin{itemize}
\item DTFT 是 DFS 当 $N\to\infty$ 的极限：两者时域离散、频域周期；但 DFS 谱离散、有谐波性、
  $X(k\Omega_0)$ 为谐波复振幅、便于计算，DTFT 谱连续、$X(\Omega)$ 为频谱密度、不便于直接计算。
\item 与 CTFT：两者时域非周期、频域为频谱密度；但 $X(\Omega)$ \textbf{周期}，$X(\omega)$ 非周期。
\end{itemize}

\subsection{傅里叶变换的离散性与周期性}

\begin{center}
\begin{tabular}{lll}
\toprule
时域信号 & 变换 & 频谱特征 \\
\midrule
连续、周期 & CFS（离散谱线） & 非周期、离散 \\
连续、非周期 & CTFT & 非周期、连续（密度函数） \\
离散、周期 & DFS & 周期、离散 \\
离散、非周期 & DTFT & 周期、连续（密度函数） \\
\bottomrule
\end{tabular}
\end{center}
四条对应规律：时域\textbf{周期}性 $\Leftrightarrow$ 频域\textbf{离散}性；时域\textbf{非周期} $\Leftrightarrow$ 频域\textbf{连续}性；
时域\textbf{离散}性 $\Leftrightarrow$ 频域\textbf{周期}性；时域\textbf{连续}性 $\Leftrightarrow$ 频域\textbf{非周期}性。

\subsection{离散傅里叶变换（DFT）}

\textbf{思路一：由 DFS 截取主值。}设 $x(n)$ 为长 $N$ 的有限长序列（$0\le n\le N-1$，称\textbf{主值序列}），
按周期 $N$ 延拓得 $x_p(n)=\sum_r x(n+rN)$，作 DFS 得 $X_p(k\Omega_0)$；只取主值区间并令
$X(k)\triangleq N X_p(k\Omega_0)$，即得\textbf{离散傅里叶变换对}
\begin{align}
X(k)&=\sum_{n=0}^{N-1}x(n)e^{-j\frac{2\pi}{N}nk},
\qquad k=0,1,\cdots,N-1,\\
x(n)&=\frac{1}{N}\sum_{k=0}^{N-1}X(k)e^{j\frac{2\pi}{N}nk},
\qquad n=0,1,\cdots,N-1.
\end{align}

\textbf{思路二：由 DTFT 频域采样。}有限长序列的 DTFT 是 $\Omega$ 的连续函数，计算机无法直接算；
在主值区间 $(0,2\pi)$ 内取 $N$ 个等间隔点 $\Omega=2\pi k/N$。若时域长 $L$、频域取 $N$ 点，
必须 $L\le N$（频域采样定理），为算法统一通常取 $L=N$，则
\begin{equation}
X(k)=X(\Omega)\Big|_{\Omega=\frac{2\pi}{N}k}
=\sum_{n=0}^{N-1}x(n)e^{-j\frac{2\pi}{N}nk}.
\end{equation}

\textbf{DFT 的本质}：有限长非周期序列原本连续的、周期的频谱密度 $X(\Omega)$，
在其数字频域主值区间（长 $2\pi$）上的\textbf{取样}，
\begin{equation}
X(k)=X(\Omega)\Big|_{\Omega=\frac{2\pi}{N}k},\qquad \Omega_0=\frac{2\pi}{N}.
\end{equation}
DFS 按傅里叶分析严格定义，DFT 则是“借用”DFS 截取主值得到的变换。

\textbf{矩阵形式}：令旋转因子 $W_N=e^{-j2\pi/N}$，则
\begin{equation}
X(k)=\sum_{n=0}^{N-1}x(n)W_N^{nk},\qquad
x(n)=\frac{1}{N}\sum_{k=0}^{N-1}X(k)W_N^{-nk},
\end{equation}
\begin{equation}
\mathbf{X}=\mathbf{W}_N^{\,nk}\mathbf{x},\qquad
\mathbf{x}=\frac{1}{N}\mathbf{W}_N^{-nk}\mathbf{X},
\end{equation}
两矩阵均 $N\times N$，且 $\mathbf{W}_N^{nk}=\big(\mathbf{W}_N^{nk}\big)^{\mathrm{T}}$（对称阵）。

\textbf{注意事项}：序列来自连续信号采样时，采样频率须大于最高截止频率的 2 倍；
对周期信号要截取\textbf{整周期}，否则产生泄漏；
矩形序列 $R_N(n)$ 的 DFT 为 $X(k)=N\delta(k)$（仅 $k=0$ 处为 $N$）。

\subsection{DFT 的性质}

设 $x(n)\xleftrightarrow{\ \mathrm{DFT}\ }X(k)$，$h(n)\xleftrightarrow{\ \mathrm{DFT}\ }H(k)$，长度均 $N$。

\begin{center}
\begin{tabular}{lll}
\toprule
性质 & 时域 & 频域 \\
\midrule
线性 & $ax_1(n)+bx_2(n)$ & $aX_1(k)+bX_2(k)$（长度须相同，短者补零） \\
圆周移位 & $\tilde{x}(n-m)_N R_N(n)$ & $W_N^{mk}X(k)$ \\
频域圆周移位 & $W_N^{-nl}x(n)$ & $X\big((k-l)\big)_N R_N(k)$ \\
时域圆周卷积 & $x(n)\circledast h(n)$ & $X(k)H(k)$ \\
频域圆周卷积（相乘） & $x(n)h(n)$ & $\dfrac{1}{N}X(k)\circledast H(k)$ \\
\bottomrule
\end{tabular}
\end{center}

\textbf{（1）圆周移位。}用取余表示周期延拓：若 $n=n_1+mN\ (0\le n_1\le N-1)$，则 $(n)_N=n_1$。
圆周移位为 $\tilde{x}(n-m)_N R_N(n)$，构造过程是“有限长序列 $\to$ 周期延拓 $\to$ 线性移位 $\to$
加矩形窗取主值”，直观上相当于把 $N$ 个样点排在圆周上旋转 $m$ 位（移出的样点从另一端转回）。

\textbf{（2）圆周卷积（圆卷积）。}
\begin{equation}
x(n)\circledast h(n)=\left[\sum_{m=0}^{N-1}x(m)h\big((n-m)\big)_N\right]R_N(n)
=\left[\sum_{m=0}^{N-1}h(m)x\big((n-m)\big)_N\right]R_N(n).
\end{equation}

\textbf{（3）与线性卷积的关系}：线性卷积右移时左端留零位，圆周卷积把移出值循环回左端，故结果一般不同。
若 $x(n)$ 长 $N$、$h(n)$ 长 $M$，线性卷积长 $N+M-1$；两者补零到长度 $L$，只要
\begin{equation}
L\ge N+M-1,
\end{equation}
则 $L$ 点圆周卷积与线性卷积\textbf{结果完全相同}。这是用 DFT/FFT 算线性卷积的依据。

\textbf{（4）奇偶虚实性}（$x(n)$ 为实序列）。由
\begin{equation}
X(k)=\sum_{n=0}^{N-1}x(n)\cos\!\left(\tfrac{2\pi}{N}nk\right)
-j\sum_{n=0}^{N-1}x(n)\sin\!\left(\tfrac{2\pi}{N}nk\right)
=X_r(k)+jX_i(k),
\end{equation}
$X_r(k)$ 为 $k$ 的偶函数、$X_i(k)$ 为奇函数，故实偶序列的 DFT 为实偶、实奇为虚奇、虚偶为虚偶、虚奇为实奇。
这里的“奇偶”是\textbf{圆周意义}下的对称（以 $N/2$ 为中心并视 $X(k)$ 为周期延拓），不是普通偶对称。

\section{快速傅里叶变换及应用}

\subsection{DFT 的计算量与 FFT 的基本思想}

$N$ 点 DFT 中每个 $X(k)$ 需 $N$ 次复乘、$N-1$ 次复加；$N$ 个 $X(k)$ 共需
\begin{equation}
N^2\ \text{次复数乘法},\qquad N(N-1)\ \text{次复数加法},
\end{equation}
复杂度 $O(N^2)$（$N=1024$ 时约 $10^6$ 次乘法），难以实时实现。
FFT 的思路：利用 $W_N$ 的周期性与对称性，把长序列 DFT 分解为短序列 DFT 再组合，
乘法次数降为 $O(N\log_2N)$。

\subsection{旋转因子 $W_N$ 的特性}

\begin{enumerate}
\item \textbf{特殊值}：$W_N^0=1$，$W_N^N=1$，$W_N^{N/2}=-1$，$W_N^{N/4}=-j$。
\item \textbf{正交性}：
\begin{equation}
\frac{1}{N}\sum_{k=0}^{N-1}W_N^{nk}\big(W_N^{mk}\big)^*
=\frac{1}{N}\sum_{k=0}^{N-1}W_N^{(n-m)k}
=\begin{cases}1, & n-m=\ell N,\\ 0, & n-m\ne\ell N.\end{cases}
\end{equation}
\item \textbf{周期性}：
\begin{equation}
W_N^{r+mN}=W_N^{r},\qquad\text{即 } W_N^{nk}=W_N^{(nk)_N},
\end{equation}
故 $W_N^{-nk}=W_N^{n(N-k)}=W_N^{k(N-n)}$；如 $N=4$ 时 $W^4=W^0$，$W^6=W^2$，$W^9=W^1$。
\item \textbf{对称性}：
\begin{equation}
W_N^{r+N/2}=-W_N^{r},\qquad\text{即 } W_N^{nk+N/2}=-W_N^{nk},
\end{equation}
如 $N=4$ 时 $W^3=-W^1$，$W^2=-W^0$。周期性加对称性使 $N=4$ 的 $W$ 矩阵只需两个独立值。
\item \textbf{可约性}：
\begin{equation}
W_N^{rn}=W_{N/r}^{n},\qquad W_{rN}^{rn}=W_N^{n}.
\end{equation}
\end{enumerate}

\subsection{基 2 按时间抽取（DIT）FFT 算法}

设 $N=2^M$，把 $x(n)$ 按序号奇偶抽取为偶序列 $x(2\ell)$ 与奇序列 $x(2\ell+1)$，$\ell=0,\cdots,N/2-1$：
\begin{align}
X(k)&=\sum_{n=0}^{N-1}x(n)W_N^{nk}
=\sum_{\ell=0}^{N/2-1}x(2\ell)W_N^{2\ell k}
+W_N^{k}\sum_{\ell=0}^{N/2-1}x(2\ell+1)W_N^{2\ell k}\nonumber\\
&=G(k)+W_N^{k}H(k),
\end{align}
其中（用可约性 $W_N^2=W_{N/2}$）
\begin{equation}
G(k)=\sum_{\ell=0}^{N/2-1}x(2\ell)W_{N/2}^{\ell k},\qquad
H(k)=\sum_{\ell=0}^{N/2-1}x(2\ell+1)W_{N/2}^{\ell k}
\end{equation}
分别是偶、奇子序列的 $N/2$ 点 DFT。因 $G$、$H$ 以 $N/2$ 为周期且 $W_N^{k+N/2}=-W_N^{k}$，得\textbf{蝶形递推}
\begin{align}
X(k)&=G(k)+W_N^{k}H(k),\\
X\!\left(k+\frac{N}{2}\right)&=G(k)-W_N^{k}H(k),
\end{align}
$k=0,1,\cdots,\frac{N}{2}-1$，前一半对应 $X(k)$ 的前 $N/2$ 点，后一半对应后 $N/2$ 点。

\textbf{继续分解}：$N/2$ 仍为偶数时，$G$、$H$ 各再分为两个 $N/4$ 点 DFT。以 $N=8$ 为例
\begin{equation}
G(k)=A(k)+W_N^{2k}B(k),\qquad
G\!\left(k+\frac{N}{4}\right)=A(k)-W_N^{2k}B(k),
\end{equation}
其中 $A,B,C,D$ 分别为 $x(4r)$、$x(4r+2)$、$x(4r+1)$、$x(4r+3)$ 的 $N/4$ 点 DFT，$H$ 同理。
$N=8$ 共三级：$1$ 个 8 点 DFT $\to$ $2$ 个 4 点 DFT $\to$ $4$ 个 2 点 DFT，级数为 $\log_2N=3$。

\textbf{蝶形运算单元}：输入两数、输出两数
\begin{equation}
\begin{aligned}
X_1&=G+W_N^{k}H,\\
X_2&=G-W_N^{k}H,
\end{aligned}
\end{equation}
每个蝶形需 1 次复乘、2 次复加；每级有 $N/2$ 个蝶形，故每级需 $N/2$ 次乘法、$N$ 次加减法，
各蝶形相互独立可并行。

\textbf{码位倒置}：按时间抽取相当于不断按奇偶分组，最终输入序列的次序是原下标的二进制码位倒置
（$N=8$ 时 $1\to4$，$2\to2$，$3\to6$……），输出 $X(k)$ 为自然顺序。

\subsection{FFT 的计算量与特点}

\begin{center}
\begin{tabular}{lll}
\toprule
算法 & 复数乘法次数 & 复数加法次数 \\
\midrule
DFT（$N$ 点） & $N^2$ & $N(N-1)$ \\
基 2 FFT（$N=2^M$） & $\dfrac{N}{2}\log_2N$ & $N\log_2N$ \\
\bottomrule
\end{tabular}
\end{center}
复杂度由 $O(N^2)$ 降为 $O(N\log_2N)$：$N=1024$ 时乘法由约 $10^6$ 降到约 $5\times10^3$，快两个数量级以上，
$N$ 越大收益越显著。

\subsection{FFT 的应用与频谱分析中的问题}

\textbf{（1）用 FFT 求线性卷积（快速卷积）}
\begin{enumerate}
\item $x(n)$ 长 $N_1$、$h(n)$ 长 $N_2$，补零到 $L\ge N_1+N_2-1$；
\item 分别作 $L$ 点 FFT 得 $X(k)$、$H(k)$，相乘得 $Y(k)=X(k)H(k)$；
\item 对 $Y(k)$ 作 $L$ 点 IFFT，得线性卷积（长 $N_1+N_2-1$）。
\end{enumerate}
数字滤波器中 $H(k)$ 可预先算好存储，实际只需两次 FFT；两序列长度相近且 $N$ 较大时，
圆周（快速）卷积比直接线性卷积省得越多；一短一长时可用分段卷积。

\textbf{（2）频率分辨率}：DFT 谱线间最小间隔，等于基波频率
\begin{equation}
f_0=\frac{1}{NT_s}=\frac{f_s}{N}\quad\left(\text{或 }\Omega_0=\frac{2\pi}{N}\right),
\end{equation}
$f_0$ 越小分辨率越高。可见\textbf{分辨率由数据长度 $NT_s$ 决定}：提高 $f_s$ 只扩大可分析频带，
要提高分辨率必须加长观测时间。

\textbf{（3）连续信号的频谱分析}
\begin{itemize}
\item \textbf{时限信号}：频带无限，$T_s$ 再小也有混叠，而过分减小 $T_s$ 又使计算量、存储量激增；
  做法是先经抗混叠滤波器滤除次要高频再采样，并选合适 $T_s$ 把混叠误差限制在允许范围内。
\item \textbf{带限信号}：时宽无限，须\textbf{加窗截断}；截断不当产生\textbf{频谱泄漏}（频谱由受限图形向外扩展）。
  矩形窗时域突变、高频衰减慢，泄漏最重；三角形窗、升余弦窗（Hanning）、改进升余弦窗（Hamming）等低旁瓣窗可减弱泄漏。
\item \textbf{连续周期信号}：也应\textbf{整周期截断}，否则产生泄漏。
\end{itemize}

\textbf{（4）幅值的标定}（易漏项）
\begin{itemize}
\item $X_s(\omega)=\frac{1}{T_s}\sum_nX(\omega-n\omega_s)$，故由 DFT 结果近似连续频谱时须\textbf{乘 $T_s$}；
  由频谱合成波形时 IDFT 结果须\textbf{乘 $f_s$}。
\item 用 DFT 求周期函数的傅里叶级数近似时，把因子 $1/N$ 放在\textbf{正变换}中更方便（正、逆变换都直接给出结果）。
\item 基 2 FFT 要求 $N$ 为 2 的整数次幂，否则补零凑成 $2^M$。
\end{itemize}

\section{离散信号的 Z 域分析}

\subsection{从 DTFT 到 Z 变换}

增长型序列（如 $a^nu(n)$，$a>1$）的 DTFT 不收敛。仿照拉普拉斯变换，把 $x(n)$ 乘衰减实指数 $r^{-n}\ (r>1)$ 后再作 DTFT：
\begin{equation}
\mathrm{DTFT}\big[x(n)r^{-n}\big]
=\sum_{n=-\infty}^{\infty}x(n)r^{-n}e^{-j\Omega n}
=\sum_{n=-\infty}^{\infty}x(n)\big(re^{j\Omega}\big)^{-n}.
\end{equation}
令 $z=re^{j\Omega}$，即得\textbf{（双边）Z 变换}
\begin{align}
X(z)&=\sum_{n=-\infty}^{\infty}x(n)z^{-n},\\
x(n)&=\frac{1}{2\pi j}\oint_{C}X(z)z^{n-1}\,\mathrm{d}z ,
\end{align}
$z^{-n}$ 的系数即 $x(n)$，幂次表示序号；围线 $C$ 是以原点为心、半径 $r$ 的逆时针闭合圆周，须位于收敛域内。
记 $\mathcal{Z}[x(n)]=X(z)$，$x(n)\xleftrightarrow{\ \mathcal{Z}\ }X(z)$。

\textbf{两个基本变换对}
\begin{align}
x(n)&=a^n u(n): & X(z)&=\frac{1}{1-az^{-1}}=\frac{z}{z-a},\quad |z|>|a|;\\
x(n)&=-a^n u(-n-1): & X(z)&=\frac{1}{1-az^{-1}}=\frac{z}{z-a},\quad |z|<|a|.
\end{align}
两式\textbf{代数表达式完全相同、收敛域不同}，故
\begin{equation}
\text{Z 变换式}+\text{收敛域}\ \Longleftrightarrow\ \text{序列}.
\end{equation}
单写 $X(z)$ 而不写收敛域无法确定对应序列——这是 Z 变换部分最重要的观念。

\subsection{收敛域与 Z 变换的几何表示}

\textbf{收敛域（ROC）}：使 $\sum x(n)z^{-n}$ 收敛的全部 $z$。常用判别法（$a_n$ 为级数第 $n$ 项）：

\begin{center}
\begin{tabular}{lll}
\toprule
判别法 & 判据 & 结论 \\
\midrule
比值判别法 & $\lim\limits_{n\to\infty}\left|\dfrac{a_{n+1}}{a_n}\right|=\rho$ & $\rho<1$ 收敛，$\rho>1$ 发散，$\rho=1$ 待定 \\
根值判别法 & $\lim\limits_{n\to\infty}\sqrt[n]{|a_n|}=\rho$ & 同上 \\
\bottomrule
\end{tabular}
\end{center}

\textbf{几何表示}：在 $z$ 平面标出零点（$\bigcirc$）与极点（$\times$）并画出收敛域。要点：
收敛域内\textbf{不含极点}；收敛域边界总以极点为界；右边序列 ROC 为某圆外部 $|z|>r_1$，
左边序列为某圆内部 $|z|<r_2$，双边序列为环域 $r_1<|z|<r_2$（此时才可能存在）。

\subsection{单边 Z 变换}

\begin{equation}
X(z)=\sum_{n=0}^{\infty}x(n)z^{-n}.
\end{equation}
单边 Z 变换可看作 $x(n)u(n)$ 的双边 Z 变换——求和只在 $n\ge0$ 上进行，不管 $n<0$ 时 $x(n)$ 是否为零；
对\textbf{因果序列}（$n<0$ 时 $x(n)=0$）二者一致；单边 Z 变换通常不特别强调收敛域（其 ROC 总为某圆外部）。

\subsection{Z 变换与其他变换的关系}

\textbf{（1）与拉普拉斯变换。}对冲激采样信号 $x_s(t)=\sum_n x(nT)\delta(t-nT)$ 取拉普拉斯变换得
$X_s(s)=\sum_n x(nT)e^{-snT}$，令 $z=e^{sT}$（即 $s=\frac{1}{T}\ln z$），得
\begin{equation}
X_s(s)\Big|_{s=\frac{1}{T}\ln z}=\sum_{n=-\infty}^{\infty}x(n)z^{-n}=X(z),
\qquad\text{即}\quad
\mathcal{L}[x_s(t)]\Big|_{s=\frac{1}{T}\ln z}=\mathcal{Z}[x(n)].
\end{equation}
即：序列的 Z 变换 = 产生该序列的理想冲激采样信号的拉普拉斯变换经 $z=e^{sT}$ 映射的结果。

\textbf{（2）$s$ 平面到 $z$ 平面的映射}。设 $s=\sigma+j\omega$、$z=re^{j\Omega}$，则
\begin{equation}
z=e^{sT}=e^{\sigma T}e^{j\omega T}
\ \Longrightarrow\
r=e^{\sigma T}=e^{\frac{2\pi}{\omega_s}\sigma},\qquad
\Omega=\omega T=\frac{2\pi\omega}{\omega_s}.
\end{equation}

\begin{center}
\begin{tabular}{ll}
\toprule
$s$ 平面 & $z$ 平面 \\
\midrule
左半平面 $\sigma<0$ & 单位圆内部 $r<1$ \\
右半平面 $\sigma>0$ & 单位圆外部 $r>1$ \\
虚轴 $\sigma=0$ & 单位圆上 $r=1$ \\
$\omega$ 每增加 $\omega_s$ & 沿单位圆逆时针转一圈 \\
\bottomrule
\end{tabular}
\end{center}
映射\textbf{不是单值的}：$s$ 平面上宽 $\omega_s$ 的每条带都映射到整个 $z$ 平面，
故虚轴上沿 $\omega$ 平移 $\omega_s$ 对应 $z$ 沿单位圆转一整圈——这正是 DTFT 频谱以 $2\pi$ 为周期的几何解释。

\textbf{（3）与 DTFT}：Z 变换是 $x(n)$ 乘 $r^{-n}$ 后的 DTFT；单位圆上的 Z 变换就是 DTFT（前提：单位圆在 ROC 内）
\begin{equation}
X(z)\Big|_{z=e^{j\Omega}}=\sum_{n=-\infty}^{\infty}x(n)e^{-j\Omega n}=X(\Omega).
\end{equation}

\textbf{（4）与 DFT}：DFT 是 Z 变换在单位圆上的 $N$ 点等间隔采样
\begin{equation}
X(z)\Big|_{z=e^{j\frac{2\pi}{N}k}}
=\sum_{n=0}^{N-1}x(n)e^{-j\frac{2\pi}{N}kn}=X(k).
\end{equation}

\section{本章小结}

\textbf{全章变换对照}：CFS（连续周期，离散非周期谱）、CTFT（连续非周期，连续非周期谱）、
DFS（离散周期，离散周期谱）、DTFT（离散非周期，连续周期谱）、DFT（时频均离散）、
Z 变换（$z$ 平面极点与收敛域）。

\textbf{主线与关键结论}
\begin{itemize}
\item \textbf{采样}：$F_s(\omega)=\sum_nP_nF(\omega-n\omega_s)$，冲激采样时 $F_s=\frac{1}{T_s}\sum_nF(\omega-n\omega_s)$；
  采样定理 $\omega_s\ge2\omega_m$；混叠不可逆，工程取 $\omega_s>10\omega_m$ 并加抗混叠滤波。
\item \textbf{时域}：$y(n)=x(n)*h(n)=\sum_mx(m)h(n-m)$，$h(n)$ 完全刻画 LTI 系统；卷积和长 $N_1+N_2-1$。
\item \textbf{频域}：DFS $\to$（$N\to\infty$）DTFT $\to$（频域采样）DFT；对偶规律“时域周期$\leftrightarrow$频域离散、
  时域离散$\leftrightarrow$频域周期”。
\item \textbf{DFT 与 FFT}：$N$ 点 DFT 为 $O(N^2)$，基 2 FFT 降为 $O(N\log_2N)$，核心是 $W_N$ 的周期性与对称性；
  用 FFT 算线性卷积须补零到 $L\ge N_1+N_2-1$。
\item \textbf{Z 域}：$X(z)=\sum_nx(n)z^{-n}$，$z=e^{sT}$；左半平面$\to$单位圆内；单位圆上为 DTFT，
  单位圆上 $N$ 点采样为 DFT；\textbf{Z 变换式必须与收敛域一起才对应唯一序列}。
\end{itemize}

% ========== 第3章 信号处理基础（依据 26春_第3章_讲义） ==========
% 讲义 3.1 系统及其性质 / 3.2 信号的线性系统处理（时域、频域、复频域）

\chapter{信号处理基础}

\section{系统及其性质}

\subsection{系统的描述}

\subsubsection{系统的定义与输入输出}
\begin{itemize}
\item \textbf{系统}：由若干相互作用、相互依赖的事物组合而成、具有特定功能的整体，可以是自然的，也可以是人造的。
\item 施加于系统的信号称为\textbf{输入信号}（激励）$x(t)$、$x(n)$，由此产生的信号称为\textbf{输出信号}（响应）$y(t)$、$y(n)$。
\item 系统的功能体现在“接受一定输入时产生什么样的输出”；任何信号的改变都要通过某种系统实现，因此\textbf{系统是信号处理的工具}（如心电信号经低通滤波去掉高频噪声）。
\end{itemize}

\subsubsection{系统的数学模型}
\begin{itemize}
\item \textbf{定义}：系统物理特性的数学抽象，以数学表达式或具有理想特性的符号组合图形表征系统特性。
\item \textbf{常见形式}：时域模型（微分/差分方程）、频域模型（频率响应）、方框图和信号流图、状态空间模型。
\item 同一物理系统在不同条件下（不同稳态工作点、不同线性化近似）可得不同形式的模型；不同物理系统经抽象近似也可能得到形式完全相同的模型。
\item 若数学模型是高阶微分方程，则该\textbf{微分方程的阶次就是系统的阶数}。
\end{itemize}

连续时间系统由线性常系数微分方程描述，离散时间系统由线性常系数差分方程描述：
\begin{align}
\sum_{k=0}^{n} a_k\, y^{(k)}(t) &= \sum_{k=0}^{m} b_k\, x^{(k)}(t), \label{eq:ctsys}\\
\sum_{k=0}^{N} a_k\, y(n-k) &= \sum_{k=0}^{M} b_k\, x(n-k). \label{eq:dtsys}
\end{align}

\subsubsection{输入输出模型与状态空间模型}
\begin{center}
\begin{tabular}{lll}
\toprule
模型 & 反映的内容 & 描述方式 \\
\midrule
输入输出模型 & 只反映系统输入与输出的关系（外特性） & 输入输出方程 \\
状态空间模型 & 既反映外特性，又着重反映系统内部状态 & 状态方程 $+$ 输出方程 \\
\bottomrule
\end{tabular}
\end{center}

\subsubsection{系统的分析与综合}
\begin{itemize}
\item \textbf{系统分析（Analysis）}：给定系统，研究系统对输入信号产生的响应，由此获得对系统功能和特性的认识，如响应求解、稳定性分析等。这类问题一般有固定的答案。
\item \textbf{系统综合（Synthesis）}：已知输入信号及对输出的要求，通过调整系统中可变动部分的结构和参数，保证得到所要求的输出，如控制器、滤波器的设计。
\end{itemize}

\subsubsection{系统的分类}
\begin{itemize}
\item \textbf{连续时间系统与离散时间系统}：前者输入输出均为连续时间信号，用微分方程或连续时间状态方程描述；后者输入输出均为离散时间信号，用差分方程或离散时间状态方程描述。
\item \textbf{单输入单输出（SISO）系统与多输入多输出（MIMO）系统}：前者单个输入、单个输出；后者有多个输入和（或）多个输出。
\item 按数学模型阶次可分为一阶、二阶、高阶系统。
\end{itemize}

\subsection{系统的性质}

\subsubsection{六条基本性质}
\begin{enumerate}
\item \textbf{记忆性}。每一时刻的输出仅取决于同时刻的输入、与过去的工作状态（历史）无关，则系统\textbf{无记忆}（瞬时系统、即时系统），由代数方程描述，如电阻 $U_R=RI_R$、加法器；否则为\textbf{有记忆}（动态系统），由微分/差分方程描述，如含储能元件的系统，$U_C(t)=\dfrac{1}{C}\displaystyle\int_{-\infty}^{t}I_C(\tau)d\tau$。
\item \textbf{因果性}。任何时刻的输出只取决于该时刻及以前的输入、与将来时刻的输入无关，则系统\textbf{因果}，否则\textbf{非因果}。电阻、电感、电容构成的实际物理系统都是因果系统；信号处理中待处理信号已被记录保存，可用后一时刻的输入决定前一时刻的输出，从而构成（允许存在的）非因果系统。
\item \textbf{可逆性}。不同输入产生不同输出（输入输出一一对应）则\textbf{可逆}，否则不可逆。如通信中的编码器应当可逆，以便接收后恢复原信号。
\item \textbf{稳定性}。对有界输入，其\textbf{零状态响应}也有界，则系统稳定，称有界输入有界输出（BIBO）稳定系统。渐近稳定是 BIBO 稳定；但 BIBO 稳定不一定满足渐近稳定条件（如传递函数存在零极点对消）。
\item \textbf{时不变性}。参数不随时间变化称\textbf{时不变}（定常）系统，参量随时间改变称\textbf{时变}（参变）系统，如 $\dot{x}=Ax$ 与 $\dot{x}=A(t)x$。
\item \textbf{线性}。同时满足\textbf{叠加性}与\textbf{齐次性}的系统为线性系统，否则为非线性系统。
\end{enumerate}

叠加性与齐次性可合并写为
\begin{equation}
x_1(t)\to y_1(t),\quad x_2(t)\to y_2(t)\ \Longrightarrow\ K_1x_1(t)+K_2x_2(t)\to K_1y_1(t)+K_2y_2(t).
\end{equation}
若系统输出与输入的关系为
\begin{equation}
y(t)=Ax(t)+y_0(t)\quad\Longrightarrow\quad \Delta y(t)=A\,\Delta x(t),
\end{equation}
即输出的增量与输入的增量之间呈线性关系，则称此类系统为\textbf{增量线性系统}。

\subsubsection{性质的判定与易错点}
\begin{center}
\begin{tabular}{p{2.6cm}p{5.2cm}p{5.6cm}}
\toprule
性质 & 判定方法 & 易错点 \\
\midrule
记忆性 & 看输出是否只与同时刻输入有关；无记忆必为代数方程 & 含电容、电感的电路一律有记忆 \\
因果性 & 检验 $y(t)=f\{x(t-\tau),\tau\ge 0\}$、$y(n)=f\{x(n-k),k\ge 0\}$ & $y(t)=\int_{-\infty}^{t}x(\tau)d\tau$ 因果；$y(t)=x(t+1)$ 非因果 \\
可逆性 & 输入输出是否一一对应 & 多对一（如平方器）不可逆 \\
稳定性 & 有界输入 $\Rightarrow$ 零状态响应有界 & 只对有界输入、只对零状态响应而言 \\
时不变性 & 令 $x_2(t)=x_1(t-t_0)$，检验是否有 $y_2(t)=y_1(t-t_0)$ & 时不变看“参数是否随时间变”，线性看“是否满足叠加与齐次”，二者独立 \\
线性 & 同时检验叠加性与齐次性 & 线性系统只有在\textbf{零初始状态}下才是线性的，否则是增量线性系统 \\
\bottomrule
\end{tabular}
\end{center}

\section{信号的线性系统处理}

线性系统处理信号有三条并行的路径：\textbf{时域法}（卷积积分/卷积和）、\textbf{频域法}（频率响应、无失真传输、理想低通滤波器）、\textbf{复频域法}（系统函数与传递函数）。三者描述同一系统，只是变换域不同。

\subsection{时域法分析}

\subsubsection{线性时不变系统的描述与直接求解的局限}
线性时不变（LTI）动态系统由式~\eqref{eq:ctsys}（连续）或式~\eqref{eq:dtsys}（离散）描述。直接求解系统方程有三点局限：
\begin{itemize}
\item 需要初始条件（系统输出、各阶导数的初值等）；
\item 纯数学方法不能表现系统响应的物理意义；
\item 输入信号或初始条件改变，必须重新求解。
\end{itemize}

\subsubsection{零输入响应与零状态响应}
\begin{itemize}
\item \textbf{零输入响应}：无外加激励，仅由起始状态（起始时刻系统储能）产生的响应，即由非零初始状态延续下来的输出。
\item \textbf{零状态响应}：系统在零初始状态下，对本次输入激励的响应。
\end{itemize}
\begin{equation}
\boxed{\ y(t)=y_{zi}(t)+y_{zs}(t)\ }\qquad\text{系统响应}=\text{零输入响应}+\text{零状态响应}.
\end{equation}
以 RC 电路为例：输入为电源电压，状态变量为电容电压 $v_0$。注意：\textbf{线性系统只有在零初始状态下才是线性的，否则是增量线性系统}。

\subsubsection{单位冲激响应与单位样值响应}
\textbf{定义}：系统在零初始条件下，对激励为单位冲激信号 $\delta(t)$ 的响应记为 $h(t)$；对激励为单位样值信号 $\delta(n)$ 的响应记为 $h(n)$。

$\delta(t)$ 与 $\delta(n)$ 的区别在于：
\begin{equation}
\delta(t)=\begin{cases}\infty, & t=0\\ 0, & t\ne 0\end{cases},\quad \int_{-\infty}^{\infty}\delta(t)\,dt=1;
\qquad
\delta(n)=\begin{cases}1, & n=0\\ 0, & n\ne 0\end{cases},\quad \sum_{n=-\infty}^{\infty}\delta(n)=1.
\end{equation}

\subsubsection{冲激响应的形式}
把 $x(t)=\delta(t)$、$y(t)=h(t)$ 代入式~\eqref{eq:ctsys}，得
\begin{equation}
\sum_{k=0}^{n}a_k h^{(k)}(t)=\sum_{k=0}^{m}b_k \delta^{(k)}(t).
\end{equation}
当 $t>0$ 时 $\delta(t)$ 及其各阶导数均为零，故 $h(t)$ 满足齐次方程 $\sum_k a_k h^{(k)}(t)=0$；又由因果性 $h(t)=0\ (t<0)$，所以 $h(t)$ 具有\textbf{齐次微分方程解的基本形式}。设齐次方程有 $n$ 个互不相同的单特征根 $\lambda_i$，则 $h(t)=\sum_{i=1}^{n}A_i e^{\lambda_i t}u(t)$。

$h(t)$ 的具体形式与 $n$、$m$ 的相对大小密切相关：
\begin{center}
\begin{tabular}{lll}
\toprule
情形 & $h(t)$ 的形式 & 说明 \\
\midrule
$n>m$ & $h(t)=\displaystyle\sum_{i=1}^{n}A_i e^{\lambda_i t}u(t)$ & 物理上可实现系统的一般形式 \\
$n=m$ & $h(t)=c\,\delta(t)+\displaystyle\sum_{i=1}^{n}A_i e^{\lambda_i t}u(t)$ & 出现冲激项 \\
$n<m$ & $h(t)=\displaystyle\sum_{j=0}^{m-n}c_j\delta^{(j)}(t)+\sum_{i=1}^{n}A_i e^{\lambda_i t}u(t)$ & 出现冲激及各阶冲激偶 \\
\bottomrule
\end{tabular}
\end{center}

离散情形同理，把 $x(n)=\delta(n)$、$y(n)=h(n)$ 代入式~\eqref{eq:dtsys}：
\begin{center}
\begin{tabular}{ll}
\toprule
情形 & $h(n)$ 的形式 \\
\midrule
$N>M$ & $h(n)=\displaystyle\sum_{i=1}^{N}A_i\lambda_i^{\,n}u(n)$ \\
$N\le M$ & $h(n)=\displaystyle\sum_{j=0}^{M-N}C_j\delta(n-j)+\sum_{i=1}^{N}A_i\lambda_i^{\,n}u(n)$ \\
\bottomrule
\end{tabular}
\end{center}

\subsubsection{求冲激响应的经典法步骤}
\begin{enumerate}
\item 列出系统的微分方程（或差分方程）；
\item 求特征方程的根；
\item 按上表写出齐次解（$h(t)$ 或 $h(n)$）的形式，待定系数；
\item 求各阶导数（离散情形为各次平移所得序列）；
\item 代回原微分（差分）方程；
\item 令等式两边各奇异函数（或各移位序列）的系数相等，解出待定系数。
\end{enumerate}

\subsubsection{卷积积分与卷积和}
\textbf{基本思想}：任意连续时间信号可分解为一系列冲激函数之和；已知 LTI 系统的单位冲激响应 $h(t)$，利用线性和时不变性就能确定系统对任意信号的响应。
\begin{align}
y(t)&=\int_{-\infty}^{\infty}x(\tau)h(t-\tau)\,d\tau=x(t)*h(t), \label{eq:conv}\\
y(n)&=\sum_{m=-\infty}^{\infty}x(m)h(n-m)=x(n)*h(n). \label{eq:convsum}
\end{align}
作为特例，$y(t)=\delta(t)*h(t)=h(t)$，$y(n)=\delta(n)*h(n)=h(n)$——即“以单位冲激激励时的零状态响应就是 $h(t)$”。

\subsubsection{阶跃响应与冲激响应的关系}
设单位阶跃响应为 $c(t)$、$c(n)$，则
\begin{align}
c(t)&=\int_{-\infty}^{\infty}h(\tau)u(t-\tau)\,d\tau, & h(t)&=c'(t);\\
c(n)&=\sum_{m=-\infty}^{\infty}h(m)u(n-m), & h(n)&=c(n)-c(n-1).
\end{align}
即\textbf{单位冲激（样值）响应是单位阶跃响应的微分（后向差分）}；反之阶跃响应是冲激响应的积分（累加）。

\subsection{频域法分析}

\subsubsection{频率响应}
频域法研究信号通过系统后在\textbf{频谱结构}上的变化，即输出响应随频率变化的规律，具有明确的物理意义。

考察单位冲激响应为 $h(t)$ 的 LTI 系统对复指数信号 $x(t)=e^{j\omega t}$ 的响应：系统对复指数信号的响应仍是\textbf{同频率}的复指数信号，只是幅值和相位发生改变，
\begin{equation}
y(t)=H(\omega)e^{j\omega t},\qquad H(\omega)=\mathcal{F}[h(t)]=\frac{Y(\omega)}{X(\omega)}=H(s)\Big|_{s=j\omega}.
\end{equation}
$H(\omega)$ 称为\textbf{频率特性（频率响应）函数}。对任意连续信号，由傅里叶反变换把 $x(t)$ 看成无穷多个不同频率复指数分量之和，再利用系统的齐次性与叠加性（或直接用傅里叶变换的时域卷积定理），得
\begin{equation}
Y(\omega)=X(\omega)H(\omega)\quad\Longleftrightarrow\quad y(t)=x(t)*h(t).
\end{equation}
写成分量形式
\begin{equation}
H(\omega)=\big|H(\omega)\big|e^{j\varphi_h(\omega)},
\end{equation}
其中 $|H(\omega)|$ 为\textbf{幅频特性}，$\varphi_h(\omega)$ 为\textbf{相频特性}。注意：严格来讲 $H(\omega)=\mathcal{F}[h(t)]$ 与 $H(\omega)=H(s)|_{s=j\omega}$ 并不一致，考虑到 $H(s)$ 在虚轴上无极点（稳定）的情况更为普遍，两种表示才一致。工程上：频率特性用极坐标图表示即\textbf{奈奎斯特图}，幅频、相频特性用对数坐标图表示即\textbf{伯德图}。

\subsubsection{离散系统的频率响应}
考察单位脉冲响应为 $h(n)$ 的 LTI 系统对复指数序列 $x(n)=e^{j\Omega n}$ 的响应，由卷积和得
\begin{equation}
H(\Omega)=\mathrm{DTFT}[h(n)]=\sum_{k=-\infty}^{\infty}h(k)e^{-j\Omega k}=\frac{Y(\Omega)}{X(\Omega)},\qquad Y(\Omega)=X(\Omega)H(\Omega).
\end{equation}
离散系统频率特性的两个要点：
\begin{itemize}
\item \textbf{周期性}：$H(\Omega)=H(\Omega+2\pi)$，是以 $2\pi$ 为周期的周期函数；若改用模拟频率 $\omega$ 表示，则周期为采样角频率 $\omega_s$。
\item \textbf{幅频与相频}：$H(\Omega)=|H(\Omega)|e^{j\varphi(\Omega)}$。
\end{itemize}

\subsubsection{无失真传输}
信号无失真传输是实现信号可靠传送与交换的基本条件：要求信号通过系统后\textbf{时域波形不变}，只允许幅度按比例放大（缩小）或在时间上有固定延迟。设原信号为 $x(t)$、频谱为 $X(\omega)$，则
\begin{equation}
y(t)=Kx(t-t_0)\quad\Longleftrightarrow\quad Y(\omega)=Ke^{-j\omega t_0}X(\omega).
\end{equation}
于是无失真传输系统的频率特性为
\begin{equation}
H(\omega)=\frac{Y(\omega)}{X(\omega)}=Ke^{-j\omega t_0}=\big|H(\omega)\big|e^{j\varphi_h(\omega)}.
\end{equation}
\textbf{无失真传输的频域条件}：
\begin{enumerate}
\item 幅频特性为与频率无关的常数，即在全部频带内具有恒定的放大倍数：$|H(\omega)|=K$；
\item 相频特性与频率呈线性关系：$\varphi_h(\omega)=-\omega t_0$（\textbf{线性相位}）。
\end{enumerate}
信号通过系统的延迟时间 $t_0$ 就是相频特性斜率的负值，工程上定义\textbf{群时延（群延迟）}
\begin{equation}
\tau=-\frac{d\varphi_h(\omega)}{d\omega}=t_0>0.
\end{equation}
对多频分量信号（如基波与二次谐波）而言，不产生相位失真的条件是各次谐波的相移必须与其频率成正比：
\begin{equation}
\frac{\varphi_1}{\omega_1}=\frac{\varphi_2}{2\omega_1}=t_0=\text{常数}.
\end{equation}
\textbf{时域条件}：$\mathcal{F}^{-1}[H(\omega)]=\mathcal{F}^{-1}[Ke^{-j\omega t_0}]=K\delta(t-t_0)$，即单位冲激响应也应是冲激信号，只是放大了 $K$ 倍、延迟了 $t_0$。

\textbf{讨论}：上述条件只是理想条件。任何系统都不可能在所有频率范围内同时具有平坦的幅频特性与线性相频特性，也不可能对单位冲激的响应是真正的冲激。实际中带有信息的物理信号都是\textbf{带限}的，只需在信号占据的频率范围内满足无失真传输条件。

\subsubsection{理想低通滤波器}
频率特性
\begin{equation}
H(\omega)=\begin{cases}1\cdot e^{-j\omega t_0}, & |\omega|<\omega_c\\[2pt] 0, & |\omega|>\omega_c\end{cases}
\end{equation}
其中 $0<\omega<\omega_c$ 为通带，$\omega>\omega_c$ 为阻带，$\omega_c$ 为截止频率。其单位冲激响应为
\begin{equation}
h(t)=\mathcal{F}^{-1}[H(\omega)]=\frac{1}{2\pi}\int_{-\omega_c}^{\omega_c}1\cdot e^{j\omega(t-t_0)}\,d\omega=\frac{\omega_c}{\pi}\Sa\big[\omega_c(t-t_0)\big].
\end{equation}
\begin{itemize}
\item 响应在 $t$ 为负值时已经出现，即 $h(t)\ne 0\ (t<0)$，\textbf{非因果系统，物理上不可实现}；
\item $h(t)$ 在时域上是无限的，这正对应“频域有限 $\Rightarrow$ 时域无限”。
\end{itemize}

\textbf{时域有限与频域无限的关系}（易错）：
\begin{itemize}
\item 信号时域有限 $\Rightarrow$ 频域无限：时域无限信号可通过加窗变成时域有限信号，对应\textbf{频域卷积}，导致频谱无限扩展；
\item 信号频域有限 $\Rightarrow$ 时域无限：相当于在频域上加窗，对应\textbf{时域卷积}，导致时域无限扩展；
\item 反之不成立：一个域上无限，\textbf{不一定}另一域上有限。如周期矩形脉冲在时域、频域上都是无限的，
\begin{equation}
F(\omega)=E\tau\omega_1\sum_{n=-\infty}^{\infty}\Sa\!\left(\frac{n\omega_1\tau}{2}\right)\delta(\omega-n\omega_1).
\end{equation}
\end{itemize}

\subsubsection{理想低通的阶跃响应}
由 $C(\omega)=\mathcal{F}[u(t)]\cdot H(\omega)=\left[\pi\delta(\omega)+\dfrac{1}{j\omega}\right]e^{-j\omega t_0}$ 取反变换，利用欧拉公式、奇函数在对称区间积分为零以及换元 $x=\omega(t-t_0)$，得
\begin{equation}
c(t)=\frac{1}{2}+\frac{1}{\pi}\operatorname{Si}\big[\omega_c(t-t_0)\big],\qquad
\operatorname{Si}(y)=\int_{0}^{y}\frac{\sin x}{x}\,dx .
\end{equation}
\begin{itemize}
\item 低通滤波器滤除了一些高频分量，不连续点在输出将被圆滑，产生渐变；
\item 阶跃响应的\textbf{上升时间}与低通滤波器（系统）的截止频率（带宽）\textbf{成反比}：
\begin{equation}
t_r=\frac{2\pi}{\omega_c}.
\end{equation}
截止频率越低，输出响应上升越缓慢。信号随时间的急剧改变意味着它包含许多高频分量。
\end{itemize}

\subsubsection{矩形脉冲经理想低通传输}
设 $e_1(t)=u(t)-u(t-\tau)$，借助叠加定理得
\begin{equation}
r_1(t)=\frac{1}{\pi}\Big\{\operatorname{Si}\big[\omega_c(t-t_0)\big]-\operatorname{Si}\big[\omega_c(t-t_0-\tau)\big]\Big\}.
\end{equation}
矩形脉冲经理想低通传输时，必须使脉宽 $\tau$ 与滤波器截止频率 $\omega_c$ 相适应（$\tau\gg\dfrac{2\pi}{\omega_c}$），才能得到大体上为矩形的响应。

\textbf{吉布斯现象}：用矩形窗函数滤取信号频谱时，时域响应的不连续点处出现上冲，上冲量约为 \textbf{9\%}。加大理想低通的带宽 $\omega_c$ 可以让更多高频分量通过、改善响应波形并减小上升时间，但\textbf{不能改变 9\% 上冲的强度}。

\subsection{复频域分析}

\subsubsection{复频域分析的意义}
\begin{itemize}
\item 在复频域中物理意义不清晰；
\item 但比频域法更方便、更有效：\textbf{求解微分方程}、\textbf{研究系统特性}、\textbf{系统的综合与设计}（自动控制理论中的传递函数即源于此）。
\item 主要内容：微分方程的复频域（拉普拉斯变换）求解、差分方程的 Z 域求解。
\end{itemize}

\subsubsection{系统函数}
\textbf{连续系统}：零初始条件下，系统输出的拉普拉斯变换与输入的拉普拉斯变换之比
\begin{equation}
H(s)=\frac{Y(s)}{X(s)}.
\end{equation}
\textbf{离散系统}：零初始条件下，系统输出的 Z 变换与输入的 Z 变换之比
\begin{equation}
H(z)=\frac{Y(z)}{X(z)}.
\end{equation}
若 $H(z)$ 的极点均位于单位圆之内，则频率响应可由系统函数得到
\begin{equation}
H(\Omega)=H(e^{j\Omega})=H(z)\Big|_{z=e^{j\Omega}};
\end{equation}
相应地，若 $H(s)$ 的极点均在左半 $s$ 平面，则
\begin{equation}
H(\omega)=H(s)\Big|_{s=j\omega}.
\end{equation}

\subsection{三域分析对照与易错点}
\begin{center}
\renewcommand{\arraystretch}{1.25}
\small
\begin{tabular}{p{1.9cm}p{3.2cm}p{3.4cm}p{5.4cm}}
\toprule
方法 & 数学工具 & 核心量 & 输入输出关系 \\
\midrule
时域法 & 卷积积分、卷积和 & $h(t)$、$h(n)$ & $y(t)=x(t)*h(t)$，$y(n)=x(n)*h(n)$ \\
频域法 & 傅里叶变换 & $H(\omega)=\mathcal{F}[h(t)]$ & $Y(\omega)=X(\omega)H(\omega)$ \\
复频域法 & 拉氏变换、Z 变换 & $H(s)=Y(s)/X(s)$，$H(z)=Y(z)/X(z)$ & 代数运算，极点决定稳定性 \\
\bottomrule
\end{tabular}
\end{center}

\begin{itemize}
\item \textbf{线性}只对零初始状态成立；有初始储能时系统是\textbf{增量线性}的，不能直接用叠加定理求全响应。
\item \textbf{因果性}看输出是否用到将来时刻的输入：$y(t)=x(t+1)$ 非因果，$y(t)=\int_{-\infty}^{t}x(\tau)d\tau$ 因果。
\item \textbf{BIBO 稳定}只约束零状态响应，且弱于渐近稳定（可能有零极点对消）。
\item $h(t)$ 的形式与 $n$、$m$ 的关系：$n>m$ 不含冲激项，$n=m$ 含 $c\delta(t)$，$n<m$ 还含冲激的各阶导数；离散情形对应 $N>M$ 与 $N\le M$。
\item \textbf{冲激响应与阶跃响应}互为微分（后向差分）与积分（累加）关系，方向不要记反。
\item 无失真传输必须\textbf{幅频为常数且相频为线性}，缺一不可；仅幅频平坦而相位非线性仍产生相位失真。
\item \textbf{理想低通}的 $h(t)$ 在 $t<0$ 已非零，物理不可实现；加大带宽只减小上升时间 $t_r=2\pi/\omega_c$，不改变 $9\%$ 上冲。
\item $\operatorname{Si}(y)$ 是正弦\textbf{积分}函数，与 $\Sa(x)=\sin x/x$ 不是同一函数，不要混用。
\item 离散系统 $H(\Omega)$ 以 $2\pi$ 为周期（模拟频率下以 $\omega_s$ 为周期），连续系统 $H(\omega)$ 无此周期性。
\end{itemize}

% ========== 第4章 滤波器（依据 26春_第4章_讲义） ==========
% 讲义 4.1 滤波器概述 / 4.2 模拟滤波器 / 4.3 数字滤波器

\chapter{滤波器}

\section{滤波器概述}

\subsection{滤波的概念及其基本原理}
\begin{itemize}
\item \textbf{传统意义}：消除或减弱干扰噪声、提取有用信号的过程（滤除噪声）。
\item \textbf{现代意义}：从原始信号中获取目标信息的过程（信号波形检测、参数估计等）。
\item \textbf{滤波的原理}：根据有用信号与噪声信号的不同特性，实现二者的有效分离，从而消除或减弱噪声、提取有用信号。
\item \textbf{滤波器}：实现滤波功能的系统，是一种具有一定传输特性的信号处理装置。从系统的角度看，滤波器就是在时域具有冲激响应 $h(t)$ 或脉冲响应 $h(n)$ 的\textbf{可实现的线性时不变系统}。
\end{itemize}

滤波的实质是“频域选择”：输入频谱 $X(\omega)$ 与滤波器的频率特性 $H(\omega)$ 相乘，得到输出频谱 $Y(\omega)=X(\omega)H(\omega)$。设计具有不同频率特性的滤波器，就得到不同的滤波效果。理想低通滤波器的矩形幅频特性物理上不可实现（非因果），实际滤波器只要求在通带内近似无失真、在阻带内充分衰减。

\subsection{滤波器的分类}

\subsubsection{经典滤波器与现代滤波器}
\begin{itemize}
\item \textbf{经典滤波器}：假定有用信号与希望去掉的信号具有\textbf{不同的频带}，通过设计合适频率特性去掉无用信号；当二者频谱相互重叠时无能为力。
\item \textbf{现代滤波器}：从含噪声的信号（数据序列）中估计出信号的某些特征或信号本身，估计结果具有更高信噪比；把信号与噪声都视为随机信号处理。
\end{itemize}

\subsubsection{按元件性质与通带范围分类}
\begin{center}
\small
\begin{tabular}{p{2.6cm}p{11.4cm}}
\toprule
分类依据 & 类别 \\
\midrule
构成元件性质 & 无源滤波器（$R$、$C$、$L$ 等无源元件）；有源滤波器（运放等有源器件，需能量源） \\
通带与阻带范围 & 低通（通低频阻高频）；高通（通高频阻低频）；带通（通过某频带、衰减其余）；带阻（抑制某频带、通过其余）；全通（某频带内各频率分量无衰减通过） \\
所处理信号性质 & 模拟滤波器（连续时间信号）；数字滤波器（离散时间信号） \\
\bottomrule
\end{tabular}
\end{center}

\subsubsection{模拟滤波器与数字滤波器的输入输出关系}
\begin{center}
\begin{tabular}{lll}
\toprule
类型 & 时域关系 & 频域关系 \\
\midrule
模拟滤波器 & $y(t)=h(t)*x(t)$ & $Y(\omega)=H(\omega)X(\omega)$ \\
数字滤波器 & $y(n)=h(n)*x(n)$ & $Y(\Omega)=H(\Omega)X(\Omega)$ \\
\bottomrule
\end{tabular}
\end{center}

\subsection{滤波器的技术指标}

\subsubsection{通带、阻带与过渡带}
理想的矩形幅频特性要求从一个频带到另一个频带突变，物理上不可实现。因此实际滤波器允许：幅频特性在通带和阻带内有一定的衰减范围并允许有起伏，且在通带与阻带之间留出\textbf{过渡带}。
\begin{itemize}
\item \textbf{通带}（通频带）：信号以很小的衰减通过的频率范围；
\item \textbf{阻带}（阻频带）：阻止信号通过的频率范围；
\item \textbf{过渡带}：通带与阻带之间的频率范围。
\end{itemize}

\subsubsection{五项主要技术指标}
\begin{center}
\renewcommand{\arraystretch}{1.3}
\small
\begin{tabular}{p{2.4cm}p{11.4cm}}
\toprule
指标 & 定义 \\
\midrule
中心频率 $\omega_0$ & 通带、阻带截止频率的几何平均值：$\omega_0=\sqrt{\omega_p\,\omega_s}$ \\
通带波动 $\Delta\alpha$ & 通带内频率特性曲线的最大峰值与谷值之差 \\
相移 $\varphi$ & 特定频率信号通过滤波器时，输入与输出端的相位之差 \\
群延迟 $\tau_g$ & 又称包络延迟，用相移对频率的变化率衡量：$\tau_g=-\dfrac{d\varphi(\omega)}{d\omega}$；实际滤波器该导数为负，故 $\tau_g$ 通常为正 \\
衰减函数 $\alpha$ & 又称衰耗特性或工作损耗，单位 dB（设 $|H(0)|$ 归一化为 1）：$\alpha=20\lg\dfrac{|H(0)|}{|H(\omega)|}=-20\lg|H(\omega)|$ \\
\bottomrule
\end{tabular}
\end{center}

工程上用两个边界点描述衰减要求：
\begin{equation}
\alpha_p=20\lg\frac{|H(0)|}{|H(\omega_p)|}=-20\lg|H(\omega_p)|\ \ (\text{通带最大衰减}),\qquad
\alpha_s=-20\lg|H(\omega_s)|\ \ (\text{阻带最小衰减}),
\end{equation}
其中 $\omega_p$ 为通带截止频率、$\omega_s$ 为阻带截止频率。设计低通滤波器时通常取幅值下降 \textbf{3 dB} 所对应的频率 $\omega_{3\text{dB}}$ 作为通带截止频率，即 $\omega_p=\omega_{3\text{dB}}$，此时 $\alpha_p=3$ dB。

\subsubsection{3 dB 与半功率点}
通常认为，当信号的功率谱密度 $|H(\omega)|^2$ 降到峰值的一半时，该频率谐波对整体的影响就可以忽略。对于 $|H(\omega)|$ 峰值为 1 的因果滤波器，半功率点对应 $|H(\omega)|^2=\dfrac{1}{2}$，即
\begin{equation}
20\lg|H(\omega)|=10\lg|H(\omega)|^2=10\lg\frac{1}{2}\approx-3\ \text{dB}.
\end{equation}
此时通带就是 $|H(\omega)|$ 从 0 dB 的峰值点下降到不小于 $-3$ dB 的频率 $\omega$ 的集合。\textbf{注意：“一半”是以功率（而非幅值）定义的}——幅值下降到 $1/\sqrt{2}\approx0.707$。

\section{模拟滤波器}

\subsection{模拟滤波器概述}
模拟滤波器是用模拟系统处理模拟信号（连续时间信号）的滤波器，是一种选择频率的装置，故又称频率选择滤波器。

\subsubsection{系统综合与设计方法}
\textbf{系统（网络）分析}是给定系统结构及其参数求响应；\textbf{系统（网络）综合}是给出激励与响应要求，求满足要求的系统结构与参数。模拟滤波器综合设计法（“插入衰减法”或“工作参数法”）分两步：
\begin{enumerate}
\item \textbf{逼近（approximation）}：根据给定频率响应特性（通带衰减 $\alpha_p$、阻带衰减 $\alpha_s$ 等指标），寻求可实现的有理函数 $H(s)$；可选若干典型 $H(s)$ 逼近理想滤波器（巴特沃思、切比雪夫等）。
\item \textbf{实现}：由选定的 $H(s)$ 实现二端口网络（滤波器）的电路结构和参数值。
\end{enumerate}
整体思路是：\textbf{频率特性的模平方 $|H(\omega)|^2$ $\to$ 系统函数 $H(s)$}。

\subsubsection{$H(s)$ 的可实现条件}
不考虑有源器件时，待求的模拟滤波器系统函数 $H(s)$ 应满足：
\begin{itemize}
\item \textbf{稳定性}：模拟滤波器应是稳定的时不变系统（存在传递函数），即 $H(s)$ 的极点分布在 $s$ \textbf{左半平面}；
\item $H(s)$ 是具有\textbf{实系数}的 $s$ 有理函数；
\item \textbf{分子多项式的阶次不大于分母多项式的阶次}，符合实际系统的因果性；
\item $h(t)=\mathcal{L}^{-1}[H(s)]$ 是 $t$ 的实函数。
\end{itemize}

\subsubsection{由模平方函数求系统函数}
由 $H(j\omega)$ 的共轭对称性 $H(j\omega)=H^*(-j\omega)$，
\begin{equation}
\big|H(j\omega)\big|^2=H(j\omega)H(-j\omega)\Big|_{j\omega\to s}=H(s)H(-s).
\end{equation}
因此 $|H(j\omega)|^2$ 的零、极点关于虚轴\textbf{成对对称分布}：一半属于 $H(s)$，另一半属于 $H(-s)$。求法为：把 $\omega^2$ 换成 $-s^2$，解出 $H(s)H(-s)$ 的全部零极点，再按可实现条件选取。
\begin{itemize}
\item \textbf{极点}：$H(s)$ 的极点必须位于 $s$ 左半平面，$H(-s)$ 的极点必须位于右半平面；
\item \textbf{零点}：取决于是否为\textbf{最小相位系统}。若为最小相位系统，$H(s)$ 的所有零点也应分布在左半平面或 $j\omega$ 轴上（这是“一定幅频特性下相移最小”的充要条件）；若非最小相位系统，零点位置与稳定性无关，可任意选取；
\item 若有零点在 $j\omega$ 轴上：由于 $H(s)$ 系数为实数，$j\omega$ 轴上的零点必须是\textbf{偶阶重零点}，此时要把轴上的零点\textbf{平分}给 $H(s)$ 与 $H(-s)$；
\item 幅度平方函数中的常数因子 $k^2$ 也要\textbf{平分}，即 $H(s)$ 的分子系数取 $k$。
\end{itemize}

\subsection{巴特沃思低通滤波器}
以巴特沃思函数作为滤波器的传递函数，它用最高阶 Taylor 级数的形式逼近理想矩形特性，是最基本的逼近函数形式之一。

\subsubsection{幅频特性（模平方函数）}
\begin{equation}
\big|H(\omega)\big|^2=\frac{1}{1+\left(\dfrac{\omega}{\omega_c}\right)^{2n}},
\end{equation}
其中 $n$ 为滤波器阶数，$\omega_c$ 为截止频率：当 $\omega=\omega_c$ 时 $|H(\omega)|^2=\dfrac{1}{2}$，即 $\omega_c$ 对应 $-3$ dB 点。

\subsubsection{主要特性}
\begin{itemize}
\item \textbf{最大平坦性}：阶数为 $n$ 时，$\omega=0$ 处前 $(2n-1)$ 阶导数均为零，原点附近非常平直，故又称“最大平坦幅值滤波器”。
\item \textbf{单调下降}：幅值在通带、阻带（含过渡带）内单调递减；$|H(0)|=1$，$|H(\infty)|=0$；无论 $n$ 取何值，$|H(\omega_c)|=0.707$，3 dB 带宽恒在 $\omega=\omega_c$。
\item 阶数 $n$ 增加：通带更平坦、阻带衰减更快、过渡带变窄，趋于理想低通；$\omega>\omega_c$ 后近似以 \textbf{$20n$ dB/倍频程}下降（代价是元件增多）。相频特性较好，通带内无起伏，接近直线。
\end{itemize}

\subsubsection{极点分布与系统函数}
由 $|H(\omega)|^2$ 令 $\omega=\dfrac{s}{j}$，得
\begin{equation}
H(s)H(-s)=\frac{1}{1+(-1)^n\left(\dfrac{s}{\omega_c}\right)^{2n}},
\end{equation}
其分母为零时
\begin{equation}
\left(\frac{s}{j\omega_c}\right)^{2n}=1\quad\Longrightarrow\quad
s_k=\omega_c\exp\!\left[\,j\!\left(\frac{2k-1}{2n}\pi+\frac{\pi}{2}\right)\right],\quad k=1,2,\dots,2n .
\end{equation}
$s_k$ 即 $H(s)$ 与 $H(-s)$ 的全部极点，其分布规律为
\begin{itemize}
\item 共 $2n$ 个极点以 $\dfrac{\pi}{n}$ 为间隔均匀分布在半径为 $\omega_c$ 的圆上，该圆称为\textbf{巴特沃思圆}；
\item 所有极点以 $j\omega$ 轴为对称轴呈对称分布，$j\omega$ 轴上没有极点；
\item $n$ 为奇数时有两个极点在实轴 $s=\pm\omega_c$ 上；$n$ 为偶数时实轴上没有极点；复数极点两两共轭对称。
\end{itemize}
取全部\textbf{左半平面}极点为 $H(s)$ 的极点（右半平面对应 $H(-s)$），分子取 $\omega_c^n$ 以保证 $|H(s)|_{s=0}=1$：
\begin{equation}
H(s)=\frac{\omega_c^n}{\displaystyle\prod_{k=1}^{n}(s-s_k)} .
\end{equation}

\subsubsection{归一化与巴特沃思多项式}
为使设计具有通用性，以 $\omega_c$ 为参考频率作归一化：分子、分母各除以 $\omega_c^n$，并令
\begin{equation}
\bar{s}=\frac{s}{\omega_c}\quad(\text{归一化复频率}).
\end{equation}
归一化后，当阶数 $n$ 为偶数时
\begin{equation}
H(\bar{s})=\frac{1}{\displaystyle\prod_{k=1}^{n/2}\left(\bar{s}^2-2\cos\!\left(\frac{2k-1}{2n}\pi+\frac{\pi}{2}\right)\bar{s}+1\right)},
\end{equation}
当阶数 $n$ 为奇数时
\begin{equation}
H(\bar{s})=\frac{1}{(\bar{s}+1)\displaystyle\prod_{k=1}^{(n-1)/2}\left(\bar{s}^2-2\cos\!\left(\frac{2k-1}{2n}\pi+\frac{\pi}{2}\right)\bar{s}+1\right)} .
\end{equation}
工程上直接查“归一化频率的各阶巴特沃思多项式”表：
\begin{center}
\renewcommand{\arraystretch}{1.2}
\small
\begin{tabular}{cl}
\toprule
$n$ & 巴特沃思多项式（分母，按 $\bar{s}$ 降幂） \\
\midrule
1 & $\bar{s}+1$ \\
2 & $\bar{s}^2+\sqrt{2}\,\bar{s}+1$ \\
3 & $\bar{s}^3+2\bar{s}^2+2\bar{s}+1$ \\
4 & $\bar{s}^4+2.613\bar{s}^3+3.414\bar{s}^2+2.613\bar{s}+1$ \\
5 & $\bar{s}^5+3.236\bar{s}^4+5.236\bar{s}^3+5.236\bar{s}^2+3.236\bar{s}+1$ \\
6 & $\bar{s}^6+3.864\bar{s}^5+7.464\bar{s}^4+9.142\bar{s}^3+7.464\bar{s}^2+3.864\bar{s}+1$ \\
7 & $\bar{s}^7+4.494\bar{s}^6+10.098\bar{s}^5+14.592\bar{s}^4+14.592\bar{s}^3+10.098\bar{s}^2+4.494\bar{s}+1$ \\
8 & $\bar{s}^8+5.153\bar{s}^7+13.137\bar{s}^6+21.846\bar{s}^5+25.688\bar{s}^4+21.846\bar{s}^3+13.137\bar{s}^2+5.153\bar{s}+1$ \\
\bottomrule
\end{tabular}
\end{center}
由 $\bar{s}$ 回到实际频率用\textbf{反归一化}：$\bar{s}=\dfrac{s}{\omega_c}$，即把表中每个 $\bar{s}$ 换成 $\dfrac{s}{\omega_c}$，或等价地把分子乘以 $\omega_c^n$、分母各系数按 $\omega_c$ 的幂次缩放。

\subsubsection{阶次的确定}
巴特沃思低通滤波器的衰减函数为
\begin{equation}
\alpha=-20\lg|H(\omega)|=10\lg\left[1+\left(\frac{\omega}{\omega_c}\right)^{2n}\right].
\end{equation}
在 $\omega=\omega_p$、$\omega=\omega_s$ 处分别为
\begin{equation}
\alpha_p=10\lg\left[1+\left(\frac{\omega_p}{\omega_c}\right)^{2n}\right],\qquad
\alpha_s=10\lg\left[1+\left(\frac{\omega_s}{\omega_c}\right)^{2n}\right].
\end{equation}
联立消去 $\omega_c$ 得阶次公式
\begin{equation}
n\ \ge\ \frac{\lg\sqrt{10^{0.1\alpha_s}-1}}{\lg\dfrac{\omega_s}{\omega_c}}\qquad(\text{取满足不等式的整数}).
\end{equation}
设计低通时通常令 $\omega_c=\omega_p=\omega_{3\text{dB}}$（此时 $\alpha_p=3$ dB），于是只需由阻带指标 $\alpha_s$、$\omega_s$ 求 $n$。

\textbf{设计步骤}（巴特沃思低通）：
\begin{enumerate}
\item 由指标确定 $\omega_c$（常取 $\omega_c=\omega_p$）并算出归一化阻带频率 $\bar{\omega}_s=\omega_s/\omega_c$；
\item 用阶次公式算出 $n$，向上取整；
\item 查表得 $n$ 阶归一化传递函数 $H(\bar{s})$；
\item 反归一化（$\bar{s}=s/\omega_c$），得实际 $H(s)$。
\end{enumerate}

\subsection{切比雪夫低通滤波器}
巴特沃思滤波器的幅频特性在通带与阻带内都随频率单调变化，滤波特性简单；但它在通带内\textbf{误差分布不均匀}，靠近频带边缘误差最大，阶数较小时阻带衰减慢、与理想特性相差较远；若要求阻带迅速下降就需提高阶数，元件增多、线路复杂。若把误差\textbf{均匀分布}在通带内，就可用较低阶数达到同样指标——这通过选择具有\textbf{等波纹}特性的逼近函数来实现。

切比雪夫滤波器由切比雪夫多项式的正交函数推导而来，采用“通带内等波动、通带外衰耗单调递增”的准则逼近理想特性：
\begin{center}
\begin{tabular}{lll}
\toprule
类型 & 通带 & 阻带 \\
\midrule
切比雪夫 \textrm{I} 型 & 等波纹 & 单调下降 \\
切比雪夫 \textrm{II} 型 & 单调 & 等波纹 \\
\bottomrule
\end{tabular}
\end{center}
与巴特沃思相比：相同阶数下切比雪夫过渡带更窄、阻带衰减更快，代价是通带内有起伏且相频特性较差（群延迟波动大）；其系统函数的计算比巴特沃思复杂得多。

\subsection{模拟滤波器的频率变换}
以归一化低通滤波器为原型，通过频率变换可得到其他类型的滤波器，讲义给出的变换有两类：
\begin{itemize}
\item \textbf{低通到高通的变换}；
\item \textbf{低通到带通的变换}。
\end{itemize}
即“先设计低通原型，再作频率变换”是模拟滤波器设计的通用套路（具体变换式由对应章节给出，此处只记思路）。

\section{数字滤波器}

\subsection{数字滤波器概述}
数字滤波器是具有一定传输特性的数字信号处理装置：输入和输出都是数字信号，借助数字器件和一定的数值计算方法，对输入信号的波形或频谱进行加工处理，去掉无用成分、保留有用成分。它既可用硬件实现，也可在计算机上用软件完成。

为便于研究由模拟到数字的映射，本节起记：模拟滤波器单位冲激响应为 $h_a(t)$、系统函数为 $H_a(s)$；数字滤波器单位脉冲响应为 $h(n)$、系统函数为 $H(z)$。

\subsubsection{数字滤波器相比模拟滤波器的优点}
\begin{itemize}
\item \textbf{精度高}：模拟器件（$R$、$L$、$C$）精度难做高，数字滤波器精度由字长决定，增加字长即可提高。
\item \textbf{可靠性高}：不受温度、湿度等环境影响，没有元器件老化、感应与杂散振荡问题。
\item \textbf{灵活性高}：通过编程可随时修改滤波特性。
\item \textbf{便于大规模集成}：设计规范、便于集成；可工作于极低频率，并容易实现模拟滤波器难以实现的\textbf{线性相位}。
\end{itemize}

\subsubsection{分类与系统函数}
按频率响应的通带特性分低通、高通、带通、带阻；按冲激响应的时间特性分无限冲激响应（IIR）与有限冲激响应（FIR）；按构成方式分递归型（如 IIR）、非递归型（如 FIR）和用 FFT 实现的。

设输入序列 $x(n)$、输出序列 $y(n)$，数字滤波器用 LTI 离散系统表示：
\begin{equation}
y(n)+\sum_{k=1}^{N}a_k y(n-k)=\sum_{k=0}^{M}b_k x(n-k),
\end{equation}
其传递函数为
\begin{equation}
H(z)=\frac{Y(z)}{X(z)}=\frac{\sum\limits_{i=0}^{M}b_i z^{-i}}{1+\sum\limits_{i=1}^{N}a_i z^{-i}}
=\frac{b_0+b_1z^{-1}+\cdots+b_Mz^{-M}}{1+a_1z^{-1}+\cdots+a_Nz^{-N}} .
\end{equation}
分子决定系统的\textbf{零点}，分母决定系统的\textbf{极点}。

\subsubsection{FIR 与 IIR 的判据}
\begin{itemize}
\item 若分母系数 $a_i=0$，则 $H(z)=\sum\limits_{i=0}^{M}b_iz^{-i}$，对应
\begin{equation}
h(n)=b_0\delta(n)+b_1\delta(n-1)+\cdots+b_M\delta(n-M),
\end{equation}
单位脉冲响应时间长度有限，最多 $(M+1)$ 项，称为\textbf{有限冲激响应（FIR）滤波器}。$H(z)$ 只有单极点 $z=0$，位于单位圆内，故 \textbf{FIR 滤波器总是稳定的}。
\item 若至少有一个 $a_i\ne0$，且分母至少存在一个根不为分子所抵消，则单位脉冲响应有无限多项，称为\textbf{无限冲激响应（IIR）滤波器}。例如 $H(z)=\dfrac{b_0}{1-z^{-1}}=b_0(1+z^{-1}+z^{-2}+\cdots)$，$|z|>1$，对应 $h(n)=b_0u(n)$。
\end{itemize}

\subsection{无限冲激响应数字滤波器}

\subsubsection{设计思路（间接法）}
\begin{itemize}
\item \textbf{直接法}：计算机辅助设计方法（不作详细讨论）。
\item \textbf{间接法}：借助模拟滤波器传递函数 $H_a(s)$ 求数字滤波器 $H(z)$。步骤为：(1) 根据技术指标确定满足要求的模拟滤波器 $H_a(s)$；(2) 寻找一种把 $s$ 平面映射到 $z$ 平面的变换关系，使 $H_a(s)$ 变为所需的 $H(z)$。
\end{itemize}
为保证数字滤波器保持模拟滤波器的特性，$s\to z$ 的映射必须满足\textbf{两个基本条件}：
\begin{enumerate}
\item $s$ 平面的复频率轴（虚轴）必须映射到 $z$ 平面的\textbf{单位圆}上，以保证频率特性相互对应；
\item $s$ 平面的\textbf{左半平面}必须映射到 $z$ 平面的\textbf{单位圆以内}，以保持稳定性。
\end{enumerate}

\subsubsection{冲激响应不变法}
\textbf{遵循准则}：使数字滤波器的单位脉冲响应等于所参照模拟滤波器单位冲激响应的采样值，
\begin{equation}
h(n)=h_a(t)\Big|_{t=nT}.
\end{equation}
\textbf{设计步骤}：
\begin{enumerate}
\item 根据技术指标确定模拟滤波器 $H_a(s)$；
\item 对 $H_a(s)$ 取拉普拉斯反变换求 $h_a(t)$；
\item 以 $T$ 为采样周期对 $h_a(t)$ 采样得 $h(n)$；
\item 求 $h(n)$ 的 Z 变换，得数字滤波器 $H(z)$。
\end{enumerate}
设 $H_a(s)$ 有 $N$ 个单极点 $p_i$，作部分分式展开
\begin{equation}
H_a(s)=\sum_{i=1}^{N}\frac{K_i}{s-p_i},\qquad K_i=(s-p_i)H_a(s)\Big|_{s=p_i},
\end{equation}
则
\begin{equation}
h_a(t)=\sum_{i=1}^{N}K_i e^{p_it}u(t)\ \xrightarrow{\ \text{采样}\ }\
h(n)=\sum_{i=1}^{N}K_i e^{p_inT}u(n)\ \xrightarrow{\ \text{Z 变换}\ }\
H(z)=\sum_{i=1}^{N}\frac{K_i}{1-e^{p_iT}z^{-1}} .
\end{equation}
即\textbf{$s$ 平面的极点 $p_i$ 映射为 $z=e^{p_iT}$ 的极点}。

\subsubsection{映射关系与频谱混叠}
由 $z=e^{sT}=e^{(\sigma+j\omega)T}=e^{\sigma T}e^{j\omega T}$ 可见：
\begin{itemize}
\item $\sigma=0$（虚轴）$\Rightarrow |z|=1$（单位圆）；$\sigma<0$（左半平面）$\Rightarrow |z|<1$（圆内）；$\sigma>0$（右半平面）$\Rightarrow |z|>1$（圆外）；
\item 当 $\sigma$ 不变、$\omega$ 改变 $\dfrac{2\pi}{T}$ 的整数倍时映射值不变，故该映射\textbf{不是单值的}：$s$ 平面沿 $j\omega$ 轴被分成一条条宽 $\dfrac{2\pi}{T}$ 的水平带，每条带都重叠映射成整个 $z$ 平面。
\end{itemize}
因此 $H_a(s)$ 与 $H(z)$ 的对应关系为
\begin{equation}
H(z)\Big|_{z=e^{sT}}=\frac{1}{T}\sum_{k=-\infty}^{\infty}H_a\!\left(s+j\frac{2\pi}{T}k\right),
\qquad
H(e^{j\Omega})=\frac{1}{T}\sum_{k=-\infty}^{\infty}H_a\!\left(j\frac{\Omega+2\pi k}{T}\right).
\end{equation}
可见冲激响应不变法存在\textbf{频谱混叠}：数字滤波器的增益与采样周期 $T$ 成反比；减小采样间隔可减小混叠，但会使数字滤波器增益过高；有时取 $h(n)=T\,h_a(nT)$ 以保证转换后增益不变。

\textbf{冲激响应不变法的特点}：模拟频率与数字频率成\textbf{线性}关系 $\Omega=\omega T$；\textbf{时域逼近}好，能很好地逼近 $h_a(t)$；但 $s\to z$ 的\textbf{多值映射}易造成频谱混叠，故只适用于低通或限带的高通、带通情况。

\subsubsection{双线性变换法}
冲激响应不变法易混叠，必须寻找一种频率特性\textbf{一一对应}的变换。双线性变换利用“微分方程表示为差分方程”的原理，把 $s$ 域 $\omega$ 从 $0$ 到 $\pm\infty$ 的分布\textbf{压缩}到 $0$ 到 $\pm\pi$ 之间，从而在 $z$ 域不出现周期重复，构成单值映射：
\begin{equation}
s=\frac{2}{T}\cdot\frac{1-z^{-1}}{1+z^{-1}},\qquad
H(z)=H_a(s)\Big|_{s=\frac{2}{T}\frac{1-z^{-1}}{1+z^{-1}}},
\end{equation}
反解得
\begin{equation}
z=\frac{1+\dfrac{T}{2}s}{1-\dfrac{T}{2}s}.
\end{equation}
\textbf{映射特性}：它是 $s$ 平面到 $z$ 平面\textbf{一一对应}（单值）的映射；把 $s$ 平面虚轴\textbf{唯一}地映射到单位圆，使 $H(z)$ 的频率响应能模仿 $H_a(s)$、\textbf{避免混叠}；把左半平面映射到单位圆内、右半平面映射到圆外，使\textbf{稳定性不变}。

\subsubsection{频率预畸变}
在单位圆上 $z=e^{j\Omega}$（对应 $s=j\omega$），代入映射式得
\begin{equation}
e^{j\Omega}=\frac{1+j\dfrac{\omega T}{2}}{1-j\dfrac{\omega T}{2}}
\quad\Longrightarrow\quad
\Omega=2\arctan\frac{\omega T}{2},\qquad
\omega=\frac{2}{T}\tan\frac{\Omega}{2}.
\end{equation}
模拟频率与数字频率之间是\textbf{非线性}关系，这会使数字滤波器与模拟滤波器在频率响应的对应关系上发生畸变（频率越高压缩越严重）。校正办法是\textbf{频率预畸变}：先对模拟滤波器的临界频率加以畸变，使其经双线性变换后正好映射为所需频率。设数字滤波器的通带、阻带截止频率为 $\Omega_p$、$\Omega_s$，则先求
\begin{equation}
\omega_p=\frac{2}{T}\tan\frac{\Omega_p}{2},\qquad
\omega_s=\frac{2}{T}\tan\frac{\Omega_s}{2},
\end{equation}
再按预畸变后的 $\omega_p$、$\omega_s$ 设计模拟滤波器，最后作双线性变换。

\textbf{双线性变换法设计步骤}：
\begin{enumerate}
\item 由数字指标 $\Omega_p$、$\Omega_s$、$\alpha_p$、$\alpha_s$ 作预畸变，得 $\omega_p$、$\omega_s$；
\item 按 $\omega_p$、$\omega_s$、$\alpha_p$、$\alpha_s$ 设计模拟低通（如巴特沃思：定阶次、查表、反归一化）；
\item 用 $s=\dfrac{2}{T}\dfrac{1-z^{-1}}{1+z^{-1}}$ 代入 $H_a(s)$，整理成 $z^{-1}$ 的有理式得 $H(z)$。
\end{enumerate}

\subsubsection{两种方法的比较}
\begin{center}
\renewcommand{\arraystretch}{1.3}
\small
\begin{tabular}{p{2.6cm}p{5.4cm}p{5.8cm}}
\toprule
比较项 & 冲激响应不变法 & 双线性变换法 \\
\midrule
映射关系 & 多值映射，$z=e^{sT}$ & 单值（一一对应）映射 \\
频率关系 & 线性：$\Omega=\omega T$ & 非线性：$\Omega=2\arctan\dfrac{\omega T}{2}$ \\
混叠 & 存在频谱混叠 & 无混叠 \\
时域逼近 & 好，$h(n)=h_a(nT)$ & 差，时域波形不保持 \\
稳定性 & 左半平面 $\to$ 圆内，稳定 & 左半平面 $\to$ 圆内，稳定 \\
适用场合 & 低通、限带的高通/带通 & 各类滤波器，需作频率预畸变 \\
\bottomrule
\end{tabular}
\end{center}

\subsection{有限冲激响应数字滤波器}

\subsubsection{IIR 与 FIR 的优缺点}
\begin{center}
\begin{tabular}{p{2.2cm}p{6.2cm}p{5.2cm}}
\toprule
类型 & 优点 & 缺点 \\
\midrule
IIR & 含零极点，可利用模拟滤波器的设计成果，计算量小、设计简单 & 零极点的存在带来稳定性问题；相位非线性 \\
FIR & 严格线性相位；滤波器总是稳定的；可利用 FFT 快速处理 & 所需阶次比 IIR 高，运算量较大 \\
\bottomrule
\end{tabular}
\end{center}

\subsubsection{设计目标}
根据要求的频率响应 $H_d(\Omega)$，找出单位脉冲响应 $h(n)$ 为\textbf{有限长}的离散时间系统，使其频率响应 $H(\Omega)$ 尽可能逼近 $H_d(\Omega)$：
\begin{equation}
h(n)\ \text{仅在}\ 0\le n\le N-1\ \text{内非零},\qquad
H(z)=\sum_{n=0}^{N-1}h(n)z^{-n},\qquad
H(\Omega)=\sum_{n=0}^{N-1}h(n)e^{-j\Omega n}.
\end{equation}
$H(z)$ 是 $z^{-1}$ 的 $N-1$ 阶多项式，其 $N-1$ 个极点全部位于 $z$ 平面原点 $z=0$ 处。

\subsubsection{窗函数法}
\begin{enumerate}
\item 由理想滤波器的频率特性 $H_d(\Omega)$ 求出理想滤波器的脉冲响应序列 $h_d(n)$（一般为无限长、非因果）；
\item 设计窗函数 $w(n)$，用它对 $h_d(n)$ 进行\textbf{截断}（加窗）得到有限长的 $h(n)$：
\begin{equation}
h(n)=h_d(n)\,w(n);
\end{equation}
\item 分析所得数字滤波器的频率特性，检验是否满足性能指标要求。
\end{enumerate}
加窗在频域相当于与窗函数频谱作卷积，故会带来过渡带加宽与通带、阻带内的起伏（吉布斯现象）：窗函数\textbf{主瓣宽度}决定过渡带宽度，\textbf{旁瓣电平}决定阻带最小衰减，需按指标折中选取。

\subsection{本章易错点}
\begin{itemize}
\item \textbf{3 dB 是半功率点}：$|H|^2$ 降为一半，即 $|H|$ 降为 $1/\sqrt{2}\approx0.707$，不是幅值降一半。
\item 巴特沃思的 3 dB 点恒在 $\omega=\omega_c$，与阶数 $n$ 无关；$n$ 只改变过渡带陡度（$20n$ dB/倍频程）。
\item 由 $|H(j\omega)|^2$ 求 $H(s)$：极点只取左半平面；零点按最小相位条件取；$j\omega$ 轴上的零点与常数因子都要\textbf{平分}给 $H(s)$ 与 $H(-s)$。
\item 阶次公式分子为 $\lg\sqrt{10^{0.1\alpha_s}-1}$（带根号），$n$ 需\textbf{向上取整}；反归一化用 $s=\bar{s}\,\omega_c$ 代入。
\item 冲激响应不变法 $\Omega=\omega T$ 为线性关系，双线性变换为\textbf{非线性}关系、必须作频率预畸变；二者都保持稳定性，但只有双线性变换无混叠。
\item FIR 的 $H(z)$ 极点在原点，故\textbf{总是稳定}的；窗函数法截断的是\textbf{理想脉冲响应} $h_d(n)$，窗的形状与长度决定过渡带与阻带衰减。
\end{itemize}


% 可选：添加参考文献/附录（如有）
% \bibliographystyle{plain}
% \bibliography{refs}

\end{document}